% DSA5208 Project 3 report. Build: scripts/make_report.sh (figures -> generated values/tables -> tectonic -> checks).
% Every number in the body comes from generated/values.tex (\val{key}) or generated/tab_*.tex, written by
% scripts/report_macros.py from the run records, the run plan and the validation output; figures come from
% scripts/figures.py. checks/requirements_check.py refuses a release with a hand-typed number, an unfilled slot,
% a figure whose time/size/core axis is not logarithmic, or a missing E row.
\documentclass[10pt,a4paper]{article}
\linespread{1.06}
\usepackage[a4paper,margin=2.2cm]{geometry}
\usepackage{lmodern}
\usepackage[T1]{fontenc}
\PassOptionsToPackage{hyphens}{url}
\usepackage{microtype,booktabs,array,tabularx,xcolor,enumitem,caption,float,amsmath,amssymb,graphicx,listings}
\usepackage[hidelinks]{hyperref}
\usepackage{titlesec}
\titleformat{\section}{\large\bfseries\sffamily}{\thesection}{0.6em}{}
\titleformat{\subsection}{\normalsize\bfseries\sffamily}{\thesubsection}{0.6em}{}
\titlespacing*{\section}{0pt}{1.4ex plus .4ex}{0.7ex}
\titlespacing*{\subsection}{0pt}{1.1ex plus .3ex}{0.5ex}
\setlist{itemsep=1pt,topsep=2pt,parsep=0pt}
\captionsetup{font=small,labelfont={bf,sf},skip=4pt}
\setlength{\parskip}{3pt}\setlength{\parindent}{0pt}
\renewcommand{\arraystretch}{1.1}
\graphicspath{{../figures/}}
\newcommand{\code}[1]{\texttt{#1}}
\newcommand{\term}[1]{\emph{#1}}% the defining use of a non-standard term
\newcommand{\req}[1]{}
\newcommand{\setval}[2]{\expandafter\def\csname pv@#1\endcsname{#2}}
\newcommand{\val}[1]{\ifcsname pv@#1\endcsname\csname pv@#1\endcsname\else\errmessage{undefined generated value #1}\fi}
\newcommand{\pending}[1]{\fcolorbox{red!60!black}{red!6}{\scriptsize\sffamily waiting for cluster data}}
\newcommand{\pendingrow}[1]{\textcolor{red!60!black}{\sffamily\footnotesize waiting for cluster data: #1}}
\newcommand{\setslot}[3]{\expandafter\def\csname sf@#1\endcsname{#2}\expandafter\def\csname sm@#1\endcsname{#3}}
\newcommand{\slotfilled}[1]{\ifcsname sf@#1\endcsname\csname sf@#1\endcsname\else 0\fi}
\newcommand{\figslot}[3]{\begin{figure}[H]\centering
  \ifnum\slotfilled{#1}=0 \fcolorbox{red!60!black}{red!4}{\parbox[c][3.6cm][c]{#2}{\centering\sffamily\small
  Figure waiting for cluster data\\[3pt]\footnotesize\csname sm@#1\endcsname}}%
  \else\includegraphics[width=#2]{#1.pdf}\fi
  \ifnum\slotfilled{#1}=2 \\[2pt]\pendingrow{partial figure, still waiting for \csname sm@#1\endcsname}\fi
  \caption{#3}\label{fig:#1}\end{figure}}
\newcommand{\interpslot}[2]{\par\smallskip\noindent\fcolorbox{orange!70!black}{orange!6}{\parbox{\dimexpr\linewidth
  -2\fboxsep-2\fboxrule}{\small\sffamily\textbf{To write once the cluster runs are in.} #2}}\par\smallskip}
% A drafted interpretation, shown with a visible mark until the user has reviewed it (the requirements check lists
% every \draftreview and does not count its E rows as filled while one remains).
\newcommand{\draftreview}[2]{\par\smallskip\noindent{\color{blue!55!black}\sffamily\scriptsize\textbf{DRAFT-FOR-USER-REVIEW}
  (#1)}\par\nopagebreak\noindent\ignorespaces#2\par\smallskip}
% A paragraph revised under the language rules (it replaces a former DRAFT-FOR-USER-REVIEW); prints nothing, and the
% requirements check lists it.
\newcommand{\revised}[1]{}
\newcounter{algo}
\newcommand{\algotitle}[2]{\refstepcounter{algo}\label{#1}\textbf{Algorithm~\thealgo: #2}}
% generated by scripts/report_macros.py; do not edit
\setval{atlas.cpn}{96}
\setval{atlas.topcores}{96}
\setval{atlas.topnodes}{1}
\setval{bytes.classic}{96}
\setval{bytes.pipelined}{128}
\setval{cal.core}{26.3}
\setval{cal.core.n}{8}
\setval{cal.core.ofpeak}{0.09}
\setval{cal.minN}{$4\times10^{6}$}
\setval{cal.node}{438}
\setval{cal.node.n}{32}
\setval{cal.node.ofpeak}{0.71}
\setval{cal.ratio}{16.7}
\setval{cal.share}{0.23}
\setval{cg.ai2}{0.16}
\setval{cg.ai3}{0.18}
\setval{cg.byteratio}{1.33}
\setval{cg.bytes.classic}{96}
\setval{cg.bytes.pipelined}{128}
\setval{cg.flops2}{15}
\setval{cg.flops3}{17}
\setval{cg.slower}{ in every strong-scaling run and in every weak-scaling run except the 2D runs with 2 and 4 ranks, where it took 0.94 and 1.00 times the classic time}
\setval{cg.st.classic}{12}
\setval{cg.st.classic.p}{3}
\setval{cg.st.classic.stencil}{3}
\setval{cg.st.classic.xr}{6}
\setval{cg.st.pipelined}{16}
\setval{cg.st.pipelined.stencil}{3}
\setval{cg.st.pipelined.upd}{13}
\setval{cl.gpu}{Vanda}
\setval{cl.gr.serial}{Vanda}
\setval{cl.gr.strong}{Vanda}
\setval{cl.gr.sweep}{Vanda}
\setval{cl.gr.variants}{Vanda}
\setval{cl.layout}{Vanda}
\setval{cl.name}{Vanda}
\setval{cl.serial}{Vanda}
\setval{cl.serialtol}{Vanda}
\setval{cl.strong}{Vanda}
\setval{cl.used}{Vanda}
\setval{cl.variants}{Vanda}
\setval{cm.2.comp}{79.7\,\%}
\setval{cm.2.halo}{2.9\,\%}
\setval{cm.2.red}{17.0\,\%}
\setval{cm.3.comp}{72.6\,\%}
\setval{cm.3.halo}{12.4\,\%}
\setval{cm.3.red}{15.0\,\%}
\setval{cpn}{72}
\setval{env.cxx}{14.3.0}
\setval{env.gpu.cxx}{26.1-0}
\setval{env.gpu.mpi}{4.1.9a1}
\setval{env.gpu.node}{GN-A40}
\setval{env.gpus}{2}
\setval{env.jobs}{15}
\setval{env.mem}{503}
\setval{env.mempersock}{252}
\setval{env.mpi}{5.0.8}
\setval{env.serialnote}{}
\setval{generated_at}{15 Nov 2026, 05:12 UTC}
\setval{geo.3d.cube.halo}{23{,}120}
\setval{geo.3d.cube.ratio}{9.8\,\%}
\setval{geo.3d.slab.halo}{329{,}672}
\setval{geo.3d.slab.ratio}{100.0\,\%}
\setval{gpu.2.cpu.t}{1.34}
\setval{gpu.2.cpu.tab}{Table~\ref{tab:layout}}
\setval{gpu.2.cpu.what}{4 ranks $\times$ 18 threads, the 2D decomposition, classic CG and the blocking halo exchange}
\setval{gpu.2.t}{1.20}
\setval{gpu.2.two.comp}{6.02}
\setval{gpu.2.two.gain}{1.96}
\setval{gpu.2.two.halo}{0.09}
\setval{gpu.2.vsnode}{1.12}
\setval{gpu.3.cpu.t}{1.46}
\setval{gpu.3.cpu.tab}{Table~\ref{tab:variants}}
\setval{gpu.3.cpu.what}{72 ranks $\times$ 1 thread, the 2D decomposition, classic CG and the blocking halo exchange}
\setval{gpu.3.t}{1.26}
\setval{gpu.3.two.comp}{6.32}
\setval{gpu.3.two.gain}{1.77}
\setval{gpu.3.two.halo}{0.77}
\setval{gpu.3.vsnode}{1.16}
\setval{gr.K}{10}
\setval{gr.a}{1}
\setval{gr.acc.com-amazon}{1.000}
\setval{gr.acc.com-dblp}{0.651}
\setval{gr.acc.com-lj}{0.986}
\setval{gr.acc.com-orkut}{0.835}
\setval{gr.acc.com-youtube}{0.886}
\setval{gr.acc.email-Eu-core}{0.758}
\setval{gr.acc.hi}{1.000 on com-Amazon}
\setval{gr.acc.lo}{0.651 on com-DBLP}
\setval{gr.amax}{100}
\setval{gr.aslope}{-0.42}
\setval{gr.avals}{0.01,0.1,1,10,100}
\setval{gr.comm.com-lj.ghosts}{163{,}302}
\setval{gr.comm.com-lj.halo}{67\,\%}
\setval{gr.comm.com-lj.nbrs}{287}
\setval{gr.comm.com-lj.red}{26\,\%}
\setval{gr.comm.com-orkut.ghosts}{792{,}468}
\setval{gr.comm.com-orkut.halo}{50\,\%}
\setval{gr.comm.com-orkut.nbrs}{143}
\setval{gr.comm.com-orkut.red}{38\,\%}
\setval{gr.cut.com-lj.bfs.nnz}{89.5\%}
\setval{gr.cut.com-lj.natural.nnz}{84.9\%}
\setval{gr.cut.com-lj.natural.rows}{82.1\%}
\setval{gr.cut.com-orkut.bfs.nnz}{91.4\%}
\setval{gr.cut.com-orkut.natural.nnz}{89.2\%}
\setval{gr.cut.com-orkut.natural.rows}{87.8\%}
\setval{gr.datalist}{email-Eu-core, com-Amazon, com-DBLP, com-Youtube, com-LiveJournal and com-Orkut}
\setval{gr.imb.com-lj.bfs.nnz}{1.01}
\setval{gr.imb.com-lj.natural.nnz}{1.00}
\setval{gr.imb.com-lj.natural.rows}{8.52}
\setval{gr.imb.com-orkut.bfs.nnz}{1.00}
\setval{gr.imb.com-orkut.natural.nnz}{1.00}
\setval{gr.imb.com-orkut.natural.rows}{3.46}
\setval{gr.keep.com-amazon.acc}{0.072}
\setval{gr.keep.com-amazon.deg}{0.7}
\setval{gr.keep.com-amazon.itmax}{21}
\setval{gr.keep.com-amazon.itmin}{1}
\setval{gr.keep.com-amazon.itratio}{12.7}
\setval{gr.keep.com-amazon.loss}{0.928}
\setval{gr.keep.com-amazon.reach}{20}
\setval{gr.keep.com-amazon.unr}{0.928}
\setval{gr.keep.com-dblp.acc}{0.088}
\setval{gr.keep.com-dblp.deg}{0.8}
\setval{gr.keep.com-dblp.itmax}{51}
\setval{gr.keep.com-dblp.itmin}{15}
\setval{gr.keep.com-dblp.itratio}{3.0}
\setval{gr.keep.com-dblp.loss}{0.563}
\setval{gr.keep.com-dblp.reach}{50}
\setval{gr.keep.com-dblp.unr}{0.685}
\setval{gr.keep.com-lj.acc}{0.719}
\setval{gr.keep.com-lj.deg}{2.2}
\setval{gr.keep.com-lj.itmax}{233}
\setval{gr.keep.com-lj.itmin}{209}
\setval{gr.keep.com-lj.itratio}{2.8}
\setval{gr.keep.com-lj.loss}{0.267}
\setval{gr.keep.com-lj.reach}{232}
\setval{gr.keep.com-lj.unr}{0.218}
\setval{gr.keep.com-orkut.acc}{0.725}
\setval{gr.keep.com-orkut.deg}{9.5}
\setval{gr.keep.com-orkut.itmax}{448}
\setval{gr.keep.com-orkut.itmin}{419}
\setval{gr.keep.com-orkut.itratio}{2.4}
\setval{gr.keep.com-orkut.loss}{0.110}
\setval{gr.keep.com-orkut.reach}{447}
\setval{gr.keep.com-orkut.unr}{0.014}
\setval{gr.keep.com-youtube.acc}{0.488}
\setval{gr.keep.com-youtube.deg}{0.7}
\setval{gr.keep.com-youtube.itmax}{263}
\setval{gr.keep.com-youtube.itmin}{236}
\setval{gr.keep.com-youtube.itratio}{2.8}
\setval{gr.keep.com-youtube.loss}{0.398}
\setval{gr.keep.com-youtube.reach}{262}
\setval{gr.keep.com-youtube.unr}{0.352}
\setval{gr.keep.email-Eu-core.acc}{0.444}
\setval{gr.keep.email-Eu-core.deg}{4.2}
\setval{gr.keep.email-Eu-core.itmax}{47}
\setval{gr.keep.email-Eu-core.itmin}{44}
\setval{gr.keep.email-Eu-core.itratio}{2.2}
\setval{gr.keep.email-Eu-core.loss}{0.314}
\setval{gr.keep.email-Eu-core.reach}{46}
\setval{gr.keep.email-Eu-core.unr}{0.230}
\setval{gr.keeplist}{0.5, 0.25 and 0.125}
\setval{gr.keepmin}{$1/8$}
\setval{gr.keepvals}{1,0.5,0.25,0.125}
\setval{gr.lift.range}{2.8 to 6.8}
\setval{gr.longrep}{10}
\setval{gr.maj.com-amazon}{0.192}
\setval{gr.maj.com-dblp}{0.227}
\setval{gr.maj.com-lj}{0.173}
\setval{gr.maj.com-orkut}{0.122}
\setval{gr.maj.com-youtube}{0.322}
\setval{gr.maj.email-Eu-core}{0.193}
\setval{gr.maxit}{100{,}000}
\setval{gr.ndata}{6}
\setval{gr.plan.cells}{194}
\setval{gr.plan.coreh}{1{,}720}
\setval{gr.rand}{0.100}
\setval{gr.rhs}{2}
\setval{gr.s}{5}
\setval{gr.ser.bfsgain}{1.75 (com-Youtube)}
\setval{gr.ser.ns.bfs}{2.24}
\setval{gr.ser.ns.natural}{2.42}
\setval{gr.ser.slope}{1.03}
\setval{gr.ser.slopebfs}{1.05}
\setval{gr.siters}{100}
\setval{gr.smax}{20}
\setval{gr.st.com-lj.classic.E}{0.06}
\setval{gr.st.com-lj.classic.Enode}{0.26}
\setval{gr.st.com-lj.classic.S}{18.6}
\setval{gr.st.com-lj.classic.T1}{103}
\setval{gr.st.com-lj.pipelined.E}{0.05}
\setval{gr.st.com-lj.pipelined.Enode}{0.27}
\setval{gr.st.com-lj.pipelined.S}{14.7}
\setval{gr.st.com-lj.pipelined.T1}{118}
\setval{gr.st.com-lj.top}{288}
\setval{gr.st.com-orkut.classic.E}{0.10}
\setval{gr.st.com-orkut.classic.Enode}{0.30}
\setval{gr.st.com-orkut.classic.S}{13.9}
\setval{gr.st.com-orkut.classic.T1}{595}
\setval{gr.st.com-orkut.pipelined.E}{0.12}
\setval{gr.st.com-orkut.pipelined.Enode}{0.29}
\setval{gr.st.com-orkut.pipelined.S}{17.8}
\setval{gr.st.com-orkut.pipelined.T1}{617}
\setval{gr.st.com-orkut.top}{144}
\setval{gr.st.note}{ The com-Orkut runs at 288 cores did not finish within the time limit of their job, so 144 cores is the largest count for that graph here.}
\setval{gr.st.top}{288}
\setval{gr.strong.counts}{1, 2, 4, 9, 18, 36, 72 and 144}
\setval{gr.strong.counts.plan}{1, 2, 4, 9, 18, 36, 72, 144, 256 and 288}
\setval{gr.strongdata}{com-LiveJournal and com-Orkut}
\setval{gr.svals}{1,2,5,10,20}
\setval{gr.sweep.ranks}{72}
\setval{gr.sweep.ranks.email}{8}
\setval{gr.theta}{0.1}
\setval{gr.tol}{$1\times10^{-8}$}
\setval{gr.val.acc}{0.758}
\setval{gr.val.bound}{$1.4\times10^{-9}$}
\setval{gr.val.bugacc}{0.587}
\setval{gr.val.bugby}{G1 and G2}
\setval{gr.val.bugspmv}{$1.1$}
\setval{gr.val.conv}{$4.6\times10^{-11}$}
\setval{gr.val.dense}{$2.3\times10^{-15}$}
\setval{gr.val.densetol}{$1\times10^{-10}$}
\setval{gr.val.iters}{10}
\setval{gr.val.layouts}{8}
\setval{gr.val.spmv}{$2.3\times10^{-15}$}
\setval{gr.val.spmvtol}{$1\times10^{-14}$}
\setval{gr.val.tinyE}{23}
\setval{gr.val.tinyV}{13}
\setval{gr.val.where}{a Vanda compute node}
\setval{gr.val.x}{$3.8\times10^{-15}$}
\setval{gr.val.xtol}{$1\times10^{-10}$}
\setval{gr.var.com-lj.best}{the natural order with equal nonzeros, the neighbourhood collective and classic CG}
\setval{gr.var.com-lj.spread}{32.4}
\setval{gr.var.com-lj.top}{288}
\setval{gr.var.com-orkut.best}{the natural order with equal nonzeros, the neighbourhood collective and classic CG}
\setval{gr.var.com-orkut.spread}{1.79}
\setval{gr.var.com-orkut.top}{72}
\setval{gr.var.cores}{72 and 288}
\setval{gr.var.note}{ For com-Orkut that count is 72 cores (Section~\ref{sec:graphstrong}).}
\setval{gs.bytes.nnz}{4}
\setval{gs.bytes.row}{16}
\setval{gs.com-amazon.bfsratio}{1.03}
\setval{gs.com-amazon.bfsratio.fast}{1.04}
\setval{gs.com-amazon.natspread}{1\,\%}
\setval{gs.com-amazon.ns}{3.20}
\setval{gs.com-dblp.bfsratio}{1.22}
\setval{gs.com-dblp.bfsratio.fast}{1.22}
\setval{gs.com-dblp.natspread}{2\,\%}
\setval{gs.com-dblp.ns}{3.27}
\setval{gs.com-lj.bfsratio}{0.96}
\setval{gs.com-lj.bfsratio.fast}{0.96}
\setval{gs.com-lj.natspread}{6\,\%}
\setval{gs.com-lj.ns}{2.11}
\setval{gs.com-orkut.bfsratio}{1.08}
\setval{gs.com-orkut.bfsratio.fast}{1.08}
\setval{gs.com-orkut.bpnnz}{5.4}
\setval{gs.com-orkut.natspread}{1\,\%}
\setval{gs.com-orkut.ns}{2.42}
\setval{gs.com-orkut.streamns}{0.20}
\setval{gs.com-youtube.bfsratio}{1.75}
\setval{gs.com-youtube.bfsratio.fast}{1.27}
\setval{gs.com-youtube.natspread}{48\,\%}
\setval{gs.com-youtube.ns}{2.79}
\setval{gs.email-Eu-core.bfsratio}{1.31}
\setval{gs.email-Eu-core.bfsratio.fast}{1.28}
\setval{gs.email-Eu-core.natspread}{42\,\%}
\setval{gs.email-Eu-core.ns}{1.56}
\setval{gs.vsgrid}{3.4}
\setval{gt.com-lj.T}{5.55}
\setval{gt.com-lj.bfs}{32.3}
\setval{gt.com-lj.bfs.a2a}{1.56}
\setval{gt.com-lj.hyb}{8$\times$9 per node}
\setval{gt.com-lj.hyb.T}{3.64}
\setval{gt.com-lj.hyb.gain}{1.52}
\setval{gt.com-lj.hyb.nbrs}{31}
\setval{gt.com-lj.hyb.vsnode}{1.48}
\setval{gt.com-lj.nbrs}{287}
\setval{gt.com-lj.pipe}{1.44}
\setval{gt.com-lj.rows}{1.25}
\setval{gt.com-lj.top}{288}
\setval{gt.com-orkut.T}{27.6}
\setval{gt.com-orkut.bfs}{1.0}
\setval{gt.com-orkut.bfs.a2a}{1.14}
\setval{gt.com-orkut.hyb}{2$\times$36 per node}
\setval{gt.com-orkut.hyb.T}{14.8}
\setval{gt.com-orkut.hyb.gain}{1.86}
\setval{gt.com-orkut.hyb.nbrs}{1}
\setval{gt.com-orkut.hyb.vsnode}{1.86}
\setval{gt.com-orkut.nbrs}{71}
\setval{gt.com-orkut.pipe}{1.05}
\setval{gt.com-orkut.rows}{1.55}
\setval{gt.com-orkut.top}{72}
\setval{gv.com-lj.node.T}{5.40}
\setval{gv.com-lj.node.a2a}{1.18}
\setval{gv.com-lj.node.bfs}{1.05}
\setval{gv.com-lj.node.cut.bfs}{87\,\%}
\setval{gv.com-lj.node.cut.hyb}{27\,\%}
\setval{gv.com-lj.node.cut.nat}{81\,\%}
\setval{gv.com-lj.node.ghosts.bfs}{520{,}130}
\setval{gv.com-lj.node.ghosts.hyb}{2{,}056{,}701}
\setval{gv.com-lj.node.ghosts.nat}{484{,}388}
\setval{gv.com-lj.node.gsum.bfs}{29{,}854{,}512}
\setval{gv.com-lj.node.gsum.hyb}{2{,}664{,}539}
\setval{gv.com-lj.node.gsum.nat}{26{,}454{,}191}
\setval{gv.com-lj.node.h8}{1.08}
\setval{gv.com-lj.node.halo}{45\,\%}
\setval{gv.com-lj.node.hybrid}{1.10}
\setval{gv.com-lj.node.hybrid.T}{4.91}
\setval{gv.com-lj.node.imb.rows}{5.85}
\setval{gv.com-lj.node.nbrs.bfs}{64}
\setval{gv.com-lj.node.nbrs.hyb}{1}
\setval{gv.com-lj.node.nbrs.nat}{71}
\setval{gv.com-lj.node.ovl}{1.11}
\setval{gv.com-lj.node.pipe}{1.14}
\setval{gv.com-lj.node.rows}{2.04}
\setval{gv.com-orkut.node.T}{27.6}
\setval{gv.com-orkut.node.a2a}{1.18}
\setval{gv.com-orkut.node.bfs}{1.01}
\setval{gv.com-orkut.node.cut.bfs}{91\,\%}
\setval{gv.com-orkut.node.cut.hyb}{33\,\%}
\setval{gv.com-orkut.node.cut.nat}{89\,\%}
\setval{gv.com-orkut.node.ghosts.bfs}{1{,}352{,}194}
\setval{gv.com-orkut.node.ghosts.hyb}{1{,}920{,}269}
\setval{gv.com-orkut.node.ghosts.nat}{1{,}210{,}552}
\setval{gv.com-orkut.node.gsum.bfs}{64{,}258{,}050}
\setval{gv.com-orkut.node.gsum.hyb}{2{,}977{,}470}
\setval{gv.com-orkut.node.gsum.nat}{66{,}607{,}218}
\setval{gv.com-orkut.node.h8}{1.73}
\setval{gv.com-orkut.node.halo}{34\,\%}
\setval{gv.com-orkut.node.hybrid}{1.86}
\setval{gv.com-orkut.node.hybrid.T}{14.8}
\setval{gv.com-orkut.node.imb.rows}{3.46}
\setval{gv.com-orkut.node.nbrs.bfs}{71}
\setval{gv.com-orkut.node.nbrs.hyb}{1}
\setval{gv.com-orkut.node.nbrs.nat}{71}
\setval{gv.com-orkut.node.ovl}{1.07}
\setval{gv.com-orkut.node.pipe}{1.05}
\setval{gv.com-orkut.node.rows}{1.55}
\setval{gw.com-amazon.a100.acc}{0.811}
\setval{gw.com-amazon.a100.it}{7}
\setval{gw.com-amazon.a100.unr}{0.189}
\setval{gw.com-amazon.amin}{0.01}
\setval{gw.com-amazon.amin.acc}{1.000}
\setval{gw.com-amazon.amin.it}{1{,}114}
\setval{gw.com-amazon.amin.unr}{0.000}
\setval{gw.com-amazon.deg}{5.5}
\setval{gw.com-amazon.dmax}{549}
\setval{gw.com-amazon.it}{113}
\setval{gw.com-amazon.itmax}{120}
\setval{gw.com-amazon.itmin}{92}
\setval{gw.com-amazon.itratio}{3.4}
\setval{gw.com-amazon.s1.acc}{1.000}
\setval{gw.com-amazon.s1.it}{107}
\setval{gw.com-amazon.s1.unr}{0.000}
\setval{gw.com-amazon.s2.acc}{1.000}
\setval{gw.com-amazon.s2.it}{110}
\setval{gw.com-amazon.s2.unr}{0.000}
\setval{gw.com-amazon.s20.acc}{1.000}
\setval{gw.com-amazon.s20.it}{118}
\setval{gw.com-amazon.s20.unr}{0.000}
\setval{gw.com-dblp.a100.acc}{0.544}
\setval{gw.com-dblp.a100.it}{13}
\setval{gw.com-dblp.a100.unr}{0.000}
\setval{gw.com-dblp.amin}{0.01}
\setval{gw.com-dblp.amin.acc}{0.680}
\setval{gw.com-dblp.amin.it}{850}
\setval{gw.com-dblp.amin.unr}{0.000}
\setval{gw.com-dblp.deg}{6.6}
\setval{gw.com-dblp.dmax}{343}
\setval{gw.com-dblp.it}{136}
\setval{gw.com-dblp.itmax}{138}
\setval{gw.com-dblp.itmin}{135}
\setval{gw.com-dblp.itratio}{5.2}
\setval{gw.com-dblp.s1.acc}{0.533}
\setval{gw.com-dblp.s1.it}{126}
\setval{gw.com-dblp.s1.unr}{0.000}
\setval{gw.com-dblp.s2.acc}{0.562}
\setval{gw.com-dblp.s2.it}{131}
\setval{gw.com-dblp.s2.unr}{0.000}
\setval{gw.com-dblp.s20.acc}{0.713}
\setval{gw.com-dblp.s20.it}{142}
\setval{gw.com-dblp.s20.unr}{0.000}
\setval{gw.com-lj.a100.acc}{0.967}
\setval{gw.com-lj.a100.it}{56}
\setval{gw.com-lj.a100.unr}{0.000}
\setval{gw.com-lj.amin}{0.1}
\setval{gw.com-lj.amin.acc}{0.985}
\setval{gw.com-lj.amin.it}{1{,}738}
\setval{gw.com-lj.amin.unr}{0.000}
\setval{gw.com-lj.deg}{17.3}
\setval{gw.com-lj.dmax}{14{,}815}
\setval{gw.com-lj.it}{624}
\setval{gw.com-lj.itmax}{666}
\setval{gw.com-lj.itmin}{602}
\setval{gw.com-lj.itratio}{3.6}
\setval{gw.com-lj.s1.acc}{0.837}
\setval{gw.com-lj.s1.it}{588}
\setval{gw.com-lj.s1.unr}{0.000}
\setval{gw.com-lj.s2.acc}{0.981}
\setval{gw.com-lj.s2.it}{607}
\setval{gw.com-lj.s2.unr}{0.000}
\setval{gw.com-lj.s20.acc}{0.990}
\setval{gw.com-lj.s20.it}{650}
\setval{gw.com-lj.s20.unr}{0.000}
\setval{gw.com-orkut.a100.acc}{0.791}
\setval{gw.com-orkut.a100.it}{115}
\setval{gw.com-orkut.a100.unr}{0.000}
\setval{gw.com-orkut.amin}{0.1}
\setval{gw.com-orkut.amin.acc}{0.833}
\setval{gw.com-orkut.amin.it}{2{,}207}
\setval{gw.com-orkut.amin.unr}{0.000}
\setval{gw.com-orkut.deg}{76.3}
\setval{gw.com-orkut.dmax}{33{,}313}
\setval{gw.com-orkut.it}{1{,}024}
\setval{gw.com-orkut.itmax}{1{,}047}
\setval{gw.com-orkut.itmin}{1{,}001}
\setval{gw.com-orkut.itratio}{4.0}
\setval{gw.com-orkut.s1.acc}{0.657}
\setval{gw.com-orkut.s1.it}{954}
\setval{gw.com-orkut.s1.unr}{0.000}
\setval{gw.com-orkut.s2.acc}{0.763}
\setval{gw.com-orkut.s2.it}{986}
\setval{gw.com-orkut.s2.unr}{0.000}
\setval{gw.com-orkut.s20.acc}{0.889}
\setval{gw.com-orkut.s20.it}{1{,}062}
\setval{gw.com-orkut.s20.unr}{0.000}
\setval{gw.com-youtube.a100.acc}{0.886}
\setval{gw.com-youtube.a100.it}{70}
\setval{gw.com-youtube.a100.unr}{0.000}
\setval{gw.com-youtube.amin}{0.01}
\setval{gw.com-youtube.amin.acc}{0.876}
\setval{gw.com-youtube.amin.it}{5{,}420}
\setval{gw.com-youtube.amin.unr}{0.000}
\setval{gw.com-youtube.deg}{5.3}
\setval{gw.com-youtube.dmax}{28{,}754}
\setval{gw.com-youtube.it}{708}
\setval{gw.com-youtube.itmax}{746}
\setval{gw.com-youtube.itmin}{688}
\setval{gw.com-youtube.itratio}{3.0}
\setval{gw.com-youtube.s1.acc}{0.857}
\setval{gw.com-youtube.s1.it}{670}
\setval{gw.com-youtube.s1.unr}{0.000}
\setval{gw.com-youtube.s2.acc}{0.895}
\setval{gw.com-youtube.s2.it}{688}
\setval{gw.com-youtube.s2.unr}{0.000}
\setval{gw.com-youtube.s20.acc}{0.898}
\setval{gw.com-youtube.s20.it}{739}
\setval{gw.com-youtube.s20.unr}{0.000}
\setval{gw.email-Eu-core.a100.acc}{0.766}
\setval{gw.email-Eu-core.a100.it}{15}
\setval{gw.email-Eu-core.a100.unr}{0.030}
\setval{gw.email-Eu-core.amin}{0.01}
\setval{gw.email-Eu-core.amin.acc}{0.738}
\setval{gw.email-Eu-core.amin.it}{184}
\setval{gw.email-Eu-core.amin.unr}{0.030}
\setval{gw.email-Eu-core.deg}{32.0}
\setval{gw.email-Eu-core.dmax}{345}
\setval{gw.email-Eu-core.it}{102}
\setval{gw.email-Eu-core.itmax}{103}
\setval{gw.email-Eu-core.itmin}{100}
\setval{gw.email-Eu-core.itratio}{3.9}
\setval{gw.email-Eu-core.s1.acc}{0.595}
\setval{gw.email-Eu-core.s1.it}{96}
\setval{gw.email-Eu-core.s1.unr}{0.028}
\setval{gw.email-Eu-core.s2.acc}{0.607}
\setval{gw.email-Eu-core.s2.it}{99}
\setval{gw.email-Eu-core.s2.unr}{0.028}
\setval{gw.email-Eu-core.s20.acc}{0.840}
\setval{gw.email-Eu-core.s20.it}{105}
\setval{gw.email-Eu-core.s20.unr}{0.034}
\setval{iters}{100}
\setval{lay.2.1.best}{4$\times$18}
\setval{lay.2.1.spread}{1.10}
\setval{lay.2.n1.pure.comp}{13.16}
\setval{lay.2.n1.pure.halo}{0.13}
\setval{lay.2.n1.pure.red}{0.23}
\setval{lay.2.n1.single.comp}{14.65}
\setval{lay.2.n1.single.halo}{0.00}
\setval{lay.2.n1.single.red}{0.00}
\setval{lay.2.n1.two.comp}{13.33}
\setval{lay.2.n1.two.halo}{0.09}
\setval{lay.2.n1.two.red}{0.03}
\setval{lay.2.top.best}{2$\times$36}
\setval{lay.2.top.pure.comp}{2.66}
\setval{lay.2.top.pure.halo}{0.10}
\setval{lay.2.top.pure.red}{0.57}
\setval{lay.2.top.single.comp}{3.89}
\setval{lay.2.top.single.halo}{0.37}
\setval{lay.2.top.single.red}{0.30}
\setval{lay.2.top.spread}{1.43}
\setval{lay.2.top.two.comp}{2.61}
\setval{lay.2.top.two.halo}{0.24}
\setval{lay.2.top.two.red}{0.33}
\setval{lay.3.1.best}{8$\times$9}
\setval{lay.3.1.spread}{1.23}
\setval{lay.3.n1.pure.comp}{13.50}
\setval{lay.3.n1.pure.halo}{0.67}
\setval{lay.3.n1.pure.red}{0.79}
\setval{lay.3.n1.single.comp}{18.14}
\setval{lay.3.n1.single.halo}{0.00}
\setval{lay.3.n1.single.red}{0.00}
\setval{lay.3.n1.two.comp}{15.12}
\setval{lay.3.n1.two.halo}{0.75}
\setval{lay.3.n1.two.red}{0.04}
\setval{lay.3.top.best}{72$\times$1}
\setval{lay.3.top.pure.comp}{2.73}
\setval{lay.3.top.pure.halo}{0.47}
\setval{lay.3.top.pure.red}{0.56}
\setval{lay.3.top.single.comp}{5.06}
\setval{lay.3.top.single.halo}{1.66}
\setval{lay.3.top.single.red}{0.47}
\setval{lay.3.top.spread}{1.91}
\setval{lay.3.top.two.comp}{2.79}
\setval{lay.3.top.two.halo}{1.46}
\setval{lay.3.top.two.red}{0.57}
\setval{lay.cpn}{72}
\setval{lay.nodes}{one node and on 4 nodes}
\setval{lay.top.sent}{ On 4 nodes the fastest layouts are 72$\times$1 in 3D and 2$\times$36 in 2D, and the slowest takes 1.91 and 1.43 times as long.}
\setval{lay.topnodes}{4}
\setval{lay.topnodes.phrase}{4 nodes}
\setval{meas.both}{ A further run at 256 cores, on 8 nodes $\times$ 32 cores shared with other jobs and with ranks not bound to cores, reached a 3D classic efficiency of 0.26.}
\setval{meas.top}{288}
\setval{meas.where}{}
\setval{node.cpn}{72}
\setval{node.cpn2}{144}
\setval{node.cps}{36}
\setval{node.cpu}{Intel Xeon Platinum 8452Y}
\setval{node.half}{36}
\setval{node.q}{18}
\setval{node.sockets}{2}
\setval{pcm.2.classic.2.imbms}{0.78}
\setval{pcm.2.classic.2.redms}{0.79}
\setval{pcm.2.classic.2node.comp}{92.1\,\%}
\setval{pcm.2.classic.2node.halo}{1.7\,\%}
\setval{pcm.2.classic.2node.imbms}{0.12}
\setval{pcm.2.classic.2node.red}{6.2\,\%}
\setval{pcm.2.classic.2node.redms}{0.41}
\setval{pcm.2.classic.node.comp}{97.1\,\%}
\setval{pcm.2.classic.node.halo}{0.9\,\%}
\setval{pcm.2.classic.node.imbms}{0.14}
\setval{pcm.2.classic.node.red}{1.7\,\%}
\setval{pcm.2.classic.node.redms}{0.23}
\setval{pcm.2.classic.top.comp}{79.7\,\%}
\setval{pcm.2.classic.top.halo}{2.9\,\%}
\setval{pcm.2.classic.top.imbms}{0.19}
\setval{pcm.2.classic.top.red}{17.0\,\%}
\setval{pcm.2.classic.top.redms}{0.57}
\setval{pcm.2.pipelined.2node.comp}{95.3\,\%}
\setval{pcm.2.pipelined.2node.halo}{1.6\,\%}
\setval{pcm.2.pipelined.2node.red}{2.9\,\%}
\setval{pcm.2.pipelined.2node.redms}{0.28}
\setval{pcm.2.pipelined.node.comp}{97.4\,\%}
\setval{pcm.2.pipelined.node.halo}{1.4\,\%}
\setval{pcm.2.pipelined.node.red}{1.1\,\%}
\setval{pcm.2.pipelined.node.redms}{0.22}
\setval{pcm.2.pipelined.top.comp}{87.4\,\%}
\setval{pcm.2.pipelined.top.halo}{2.5\,\%}
\setval{pcm.2.pipelined.top.red}{7.9\,\%}
\setval{pcm.2.pipelined.top.redms}{0.37}
\setval{pcm.3.classic.2.imbms}{0.27}
\setval{pcm.3.classic.2.redms}{0.85}
\setval{pcm.3.classic.2node.comp}{82.8\,\%}
\setval{pcm.3.classic.2node.halo}{6.6\,\%}
\setval{pcm.3.classic.2node.imbms}{0.56}
\setval{pcm.3.classic.2node.red}{10.4\,\%}
\setval{pcm.3.classic.2node.redms}{0.79}
\setval{pcm.3.classic.node.comp}{90.2\,\%}
\setval{pcm.3.classic.node.halo}{4.5\,\%}
\setval{pcm.3.classic.node.imbms}{0.92}
\setval{pcm.3.classic.node.red}{5.3\,\%}
\setval{pcm.3.classic.node.redms}{0.79}
\setval{pcm.3.classic.top.comp}{72.6\,\%}
\setval{pcm.3.classic.top.halo}{12.4\,\%}
\setval{pcm.3.classic.top.imbms}{0.25}
\setval{pcm.3.classic.top.red}{15.0\,\%}
\setval{pcm.3.classic.top.redms}{0.56}
\setval{pcm.3.pipelined.2node.comp}{87.7\,\%}
\setval{pcm.3.pipelined.2node.halo}{7.1\,\%}
\setval{pcm.3.pipelined.2node.red}{5.1\,\%}
\setval{pcm.3.pipelined.2node.redms}{0.53}
\setval{pcm.3.pipelined.node.comp}{89.6\,\%}
\setval{pcm.3.pipelined.node.halo}{6.5\,\%}
\setval{pcm.3.pipelined.node.red}{3.5\,\%}
\setval{pcm.3.pipelined.node.redms}{0.75}
\setval{pcm.3.pipelined.top.comp}{80.1\,\%}
\setval{pcm.3.pipelined.top.halo}{7.8\,\%}
\setval{pcm.3.pipelined.top.red}{12.3\,\%}
\setval{pcm.3.pipelined.top.redms}{0.67}
\setval{pgs.com-lj.classic.2node.E}{0.09}
\setval{pgs.com-lj.classic.2node.S}{12.8}
\setval{pgs.com-lj.classic.2node.T}{8.07}
\setval{pgs.com-lj.classic.2node.halo}{52\,\%}
\setval{pgs.com-lj.classic.2node.red}{38\,\%}
\setval{pgs.com-lj.classic.node.E}{0.26}
\setval{pgs.com-lj.classic.node.S}{19.1}
\setval{pgs.com-lj.classic.node.T}{5.40}
\setval{pgs.com-lj.classic.node.halo}{45\,\%}
\setval{pgs.com-lj.classic.node.red}{21\,\%}
\setval{pgs.com-lj.classic.slowdown}{1.49}
\setval{pgs.com-lj.classic.top.E}{0.06}
\setval{pgs.com-lj.classic.top.S}{18.6}
\setval{pgs.com-lj.classic.top.T}{5.55}
\setval{pgs.com-lj.classic.top.halo}{67\,\%}
\setval{pgs.com-lj.classic.top.red}{26\,\%}
\setval{pgs.com-lj.itratio}{1.08}
\setval{pgs.com-lj.pipelined.2node.E}{0.14}
\setval{pgs.com-lj.pipelined.2node.S}{20.3}
\setval{pgs.com-lj.pipelined.2node.T}{5.80}
\setval{pgs.com-lj.pipelined.2node.halo}{72\,\%}
\setval{pgs.com-lj.pipelined.2node.red}{9\,\%}
\setval{pgs.com-lj.pipelined.node.E}{0.27}
\setval{pgs.com-lj.pipelined.node.S}{19.1}
\setval{pgs.com-lj.pipelined.node.T}{6.18}
\setval{pgs.com-lj.pipelined.node.halo}{56\,\%}
\setval{pgs.com-lj.pipelined.node.red}{8\,\%}
\setval{pgs.com-lj.pipelined.slowdown}{0.94}
\setval{pgs.com-lj.pipelined.top.E}{0.05}
\setval{pgs.com-lj.pipelined.top.S}{14.7}
\setval{pgs.com-lj.pipelined.top.T}{8.01}
\setval{pgs.com-lj.pipelined.top.halo}{92\,\%}
\setval{pgs.com-lj.pipelined.top.red}{2\,\%}
\setval{pgs.com-lj.pratio.2node}{0.72}
\setval{pgs.com-lj.pratio.node}{1.14}
\setval{pgs.com-lj.pratio.top}{1.44}
\setval{pgs.com-orkut.classic.2node.E}{0.10}
\setval{pgs.com-orkut.classic.2node.S}{13.9}
\setval{pgs.com-orkut.classic.2node.T}{42.7}
\setval{pgs.com-orkut.classic.2node.halo}{50\,\%}
\setval{pgs.com-orkut.classic.2node.red}{38\,\%}
\setval{pgs.com-orkut.classic.node.E}{0.30}
\setval{pgs.com-orkut.classic.node.S}{21.6}
\setval{pgs.com-orkut.classic.node.T}{27.6}
\setval{pgs.com-orkut.classic.node.halo}{34\,\%}
\setval{pgs.com-orkut.classic.node.red}{16\,\%}
\setval{pgs.com-orkut.classic.slowdown}{1.55}
\setval{pgs.com-orkut.classic.top.E}{\pending{pgs.com-orkut.classic.top.E}}
\setval{pgs.com-orkut.classic.top.S}{\pending{pgs.com-orkut.classic.top.S}}
\setval{pgs.com-orkut.classic.top.T}{\pending{pgs.com-orkut.classic.top.T}}
\setval{pgs.com-orkut.classic.top.halo}{\pending{pgs.com-orkut.classic.top.halo}}
\setval{pgs.com-orkut.classic.top.red}{\pending{pgs.com-orkut.classic.top.red}}
\setval{pgs.com-orkut.itratio}{1.03}
\setval{pgs.com-orkut.pipelined.2node.E}{0.12}
\setval{pgs.com-orkut.pipelined.2node.S}{17.8}
\setval{pgs.com-orkut.pipelined.2node.T}{34.6}
\setval{pgs.com-orkut.pipelined.2node.halo}{76\,\%}
\setval{pgs.com-orkut.pipelined.2node.red}{5\,\%}
\setval{pgs.com-orkut.pipelined.node.E}{0.29}
\setval{pgs.com-orkut.pipelined.node.S}{21.2}
\setval{pgs.com-orkut.pipelined.node.T}{29.1}
\setval{pgs.com-orkut.pipelined.node.halo}{41\,\%}
\setval{pgs.com-orkut.pipelined.node.red}{8\,\%}
\setval{pgs.com-orkut.pipelined.slowdown}{1.19}
\setval{pgs.com-orkut.pipelined.top.E}{\pending{pgs.com-orkut.pipelined.top.E}}
\setval{pgs.com-orkut.pipelined.top.S}{\pending{pgs.com-orkut.pipelined.top.S}}
\setval{pgs.com-orkut.pipelined.top.T}{\pending{pgs.com-orkut.pipelined.top.T}}
\setval{pgs.com-orkut.pipelined.top.halo}{\pending{pgs.com-orkut.pipelined.top.halo}}
\setval{pgs.com-orkut.pipelined.top.red}{\pending{pgs.com-orkut.pipelined.top.red}}
\setval{pgs.com-orkut.pratio.2node}{0.81}
\setval{pgs.com-orkut.pratio.node}{1.05}
\setval{pgs.com-orkut.pratio.top}{\pending{pgs.com-orkut.pratio.top}}
\setval{plan.atlas.cells}{296}
\setval{plan.cells}{350}
\setval{plan.jobs}{18}
\setval{plan.layouts}{1$\times$72, 2$\times$36, 4$\times$18, 8$\times$9 and 72$\times$1}
\setval{plan.vanda.cells}{350}
\setval{plan.variantsP}{72 and 288}
\setval{pst.2.classic.1.E}{1.00}
\setval{pst.2.classic.1.S}{1.0}
\setval{pst.2.classic.1.T}{24.1}
\setval{pst.2.classic.144.E}{0.25}
\setval{pst.2.classic.144.S}{36.3}
\setval{pst.2.classic.144.T}{0.664}
\setval{pst.2.classic.18.E}{0.43}
\setval{pst.2.classic.18.S}{7.7}
\setval{pst.2.classic.18.T}{3.14}
\setval{pst.2.classic.2.E}{1.00}
\setval{pst.2.classic.2.S}{2.0}
\setval{pst.2.classic.2.T}{12.1}
\setval{pst.2.classic.2.ns}{3.62}
\setval{pst.2.classic.288.E}{0.25}
\setval{pst.2.classic.288.S}{72.2}
\setval{pst.2.classic.288.T}{0.334}
\setval{pst.2.classic.2node.E}{0.25}
\setval{pst.2.classic.2node.S}{36.3}
\setval{pst.2.classic.2node.T}{0.664}
\setval{pst.2.classic.2node.cbw}{527}
\setval{pst.2.classic.2node.cns}{13.1}
\setval{pst.2.classic.2node.cpk}{0.86}
\setval{pst.2.classic.2node.vec}{3.73}
\setval{pst.2.classic.36.E}{0.24}
\setval{pst.2.classic.36.S}{8.8}
\setval{pst.2.classic.36.T}{2.75}
\setval{pst.2.classic.4.E}{0.91}
\setval{pst.2.classic.4.S}{3.6}
\setval{pst.2.classic.4.T}{6.62}
\setval{pst.2.classic.72.E}{0.25}
\setval{pst.2.classic.72.S}{17.7}
\setval{pst.2.classic.72.T}{1.36}
\setval{pst.2.classic.9.E}{0.64}
\setval{pst.2.classic.9.S}{5.7}
\setval{pst.2.classic.9.T}{4.20}
\setval{pst.2.classic.cgain.cores}{4}
\setval{pst.2.classic.cgain.node.top}{4.94}
\setval{pst.2.classic.gain.2node.top}{1.99}
\setval{pst.2.classic.gain.half.node}{2.02}
\setval{pst.2.classic.gain.node.2node}{2.05}
\setval{pst.2.classic.gain.q.half}{1.14}
\setval{pst.2.classic.half.E}{0.24}
\setval{pst.2.classic.half.S}{8.8}
\setval{pst.2.classic.half.T}{2.75}
\setval{pst.2.classic.node.E}{0.25}
\setval{pst.2.classic.node.S}{17.7}
\setval{pst.2.classic.node.T}{1.36}
\setval{pst.2.classic.node.cbw}{489}
\setval{pst.2.classic.node.cns}{14.1}
\setval{pst.2.classic.node.cpk}{0.80}
\setval{pst.2.classic.node.vec}{7.46}
\setval{pst.2.classic.one.E}{1.00}
\setval{pst.2.classic.one.S}{1.0}
\setval{pst.2.classic.one.T}{24.1}
\setval{pst.2.classic.one.ns}{3.59}
\setval{pst.2.classic.q.E}{0.43}
\setval{pst.2.classic.q.S}{7.7}
\setval{pst.2.classic.q.T}{3.14}
\setval{pst.2.classic.top.E}{0.25}
\setval{pst.2.classic.top.S}{72.2}
\setval{pst.2.classic.top.T}{0.334}
\setval{pst.2.classic.top.cbw}{604}
\setval{pst.2.classic.top.cns}{11.4}
\setval{pst.2.classic.top.cpk}{0.98}
\setval{pst.2.classic.top.vec}{1.86}
\setval{pst.2.pipelined.1.E}{1.00}
\setval{pst.2.pipelined.1.S}{1.0}
\setval{pst.2.pipelined.1.T}{26.5}
\setval{pst.2.pipelined.144.E}{0.20}
\setval{pst.2.pipelined.144.S}{28.3}
\setval{pst.2.pipelined.144.T}{0.936}
\setval{pst.2.pipelined.18.E}{0.36}
\setval{pst.2.pipelined.18.S}{6.4}
\setval{pst.2.pipelined.18.T}{4.12}
\setval{pst.2.pipelined.2.E}{0.99}
\setval{pst.2.pipelined.2.S}{2.0}
\setval{pst.2.pipelined.2.T}{13.4}
\setval{pst.2.pipelined.288.E}{0.20}
\setval{pst.2.pipelined.288.S}{56.7}
\setval{pst.2.pipelined.288.T}{0.467}
\setval{pst.2.pipelined.2node.E}{0.20}
\setval{pst.2.pipelined.2node.S}{28.3}
\setval{pst.2.pipelined.2node.T}{0.936}
\setval{pst.2.pipelined.36.E}{0.19}
\setval{pst.2.pipelined.36.S}{7.0}
\setval{pst.2.pipelined.36.T}{3.79}
\setval{pst.2.pipelined.4.E}{0.90}
\setval{pst.2.pipelined.4.S}{3.6}
\setval{pst.2.pipelined.4.T}{7.40}
\setval{pst.2.pipelined.72.E}{0.19}
\setval{pst.2.pipelined.72.S}{14.0}
\setval{pst.2.pipelined.72.T}{1.89}
\setval{pst.2.pipelined.9.E}{0.59}
\setval{pst.2.pipelined.9.S}{5.4}
\setval{pst.2.pipelined.9.T}{4.95}
\setval{pst.2.pipelined.half.E}{0.19}
\setval{pst.2.pipelined.half.S}{7.0}
\setval{pst.2.pipelined.half.T}{3.79}
\setval{pst.2.pipelined.node.E}{0.19}
\setval{pst.2.pipelined.node.S}{14.0}
\setval{pst.2.pipelined.node.T}{1.89}
\setval{pst.2.pipelined.one.E}{1.00}
\setval{pst.2.pipelined.one.S}{1.0}
\setval{pst.2.pipelined.one.T}{26.5}
\setval{pst.2.pipelined.q.E}{0.36}
\setval{pst.2.pipelined.q.S}{6.4}
\setval{pst.2.pipelined.q.T}{4.12}
\setval{pst.2.pipelined.top.E}{0.20}
\setval{pst.2.pipelined.top.S}{56.7}
\setval{pst.2.pipelined.top.T}{0.467}
\setval{pst.2.pratio.1}{1.10}
\setval{pst.2.pratio.144}{1.41}
\setval{pst.2.pratio.18}{1.31}
\setval{pst.2.pratio.2}{1.11}
\setval{pst.2.pratio.288}{1.40}
\setval{pst.2.pratio.2node}{1.41}
\setval{pst.2.pratio.36}{1.38}
\setval{pst.2.pratio.4}{1.12}
\setval{pst.2.pratio.72}{1.39}
\setval{pst.2.pratio.9}{1.18}
\setval{pst.2.pratio.half}{1.38}
\setval{pst.2.pratio.node}{1.39}
\setval{pst.2.pratio.one}{1.10}
\setval{pst.2.pratio.q}{1.31}
\setval{pst.2.pratio.top}{1.40}
\setval{pst.3.classic.1.E}{1.00}
\setval{pst.3.classic.1.S}{1.0}
\setval{pst.3.classic.1.T}{28.4}
\setval{pst.3.classic.144.E}{0.26}
\setval{pst.3.classic.144.S}{37.3}
\setval{pst.3.classic.144.T}{0.762}
\setval{pst.3.classic.18.E}{0.46}
\setval{pst.3.classic.18.S}{8.3}
\setval{pst.3.classic.18.T}{3.43}
\setval{pst.3.classic.2.E}{1.00}
\setval{pst.3.classic.2.S}{2.0}
\setval{pst.3.classic.2.T}{14.2}
\setval{pst.3.classic.2.ns}{4.25}
\setval{pst.3.classic.288.E}{0.26}
\setval{pst.3.classic.288.S}{75.3}
\setval{pst.3.classic.288.T}{0.377}
\setval{pst.3.classic.2node.E}{0.26}
\setval{pst.3.classic.2node.S}{37.3}
\setval{pst.3.classic.2node.T}{0.762}
\setval{pst.3.classic.2node.cbw}{509}
\setval{pst.3.classic.2node.cns}{13.6}
\setval{pst.3.classic.2node.cpk}{0.83}
\setval{pst.3.classic.2node.vec}{3.72}
\setval{pst.3.classic.36.E}{0.26}
\setval{pst.3.classic.36.S}{9.4}
\setval{pst.3.classic.36.T}{3.01}
\setval{pst.3.classic.4.E}{0.97}
\setval{pst.3.classic.4.S}{3.9}
\setval{pst.3.classic.4.T}{7.35}
\setval{pst.3.classic.72.E}{0.26}
\setval{pst.3.classic.72.S}{18.9}
\setval{pst.3.classic.72.T}{1.50}
\setval{pst.3.classic.9.E}{0.72}
\setval{pst.3.classic.9.S}{6.4}
\setval{pst.3.classic.9.T}{4.41}
\setval{pst.3.classic.cgain.cores}{4}
\setval{pst.3.classic.cgain.node.top}{4.94}
\setval{pst.3.classic.gain.2node.top}{2.02}
\setval{pst.3.classic.gain.half.node}{2.01}
\setval{pst.3.classic.gain.node.2node}{1.97}
\setval{pst.3.classic.gain.q.half}{1.14}
\setval{pst.3.classic.half.E}{0.26}
\setval{pst.3.classic.half.S}{9.4}
\setval{pst.3.classic.half.T}{3.01}
\setval{pst.3.classic.node.E}{0.26}
\setval{pst.3.classic.node.S}{18.9}
\setval{pst.3.classic.node.T}{1.50}
\setval{pst.3.classic.node.cbw}{476}
\setval{pst.3.classic.node.cns}{14.5}
\setval{pst.3.classic.node.cpk}{0.77}
\setval{pst.3.classic.node.vec}{7.44}
\setval{pst.3.classic.one.E}{1.00}
\setval{pst.3.classic.one.S}{1.0}
\setval{pst.3.classic.one.T}{28.4}
\setval{pst.3.classic.one.ns}{4.24}
\setval{pst.3.classic.q.E}{0.46}
\setval{pst.3.classic.q.S}{8.3}
\setval{pst.3.classic.q.T}{3.43}
\setval{pst.3.classic.top.E}{0.26}
\setval{pst.3.classic.top.S}{75.3}
\setval{pst.3.classic.top.T}{0.377}
\setval{pst.3.classic.top.cbw}{588}
\setval{pst.3.classic.top.cns}{11.8}
\setval{pst.3.classic.top.cpk}{0.96}
\setval{pst.3.classic.top.vec}{1.86}
\setval{pst.3.pipelined.1.E}{1.00}
\setval{pst.3.pipelined.1.S}{1.0}
\setval{pst.3.pipelined.1.T}{30.5}
\setval{pst.3.pipelined.144.E}{0.20}
\setval{pst.3.pipelined.144.S}{29.0}
\setval{pst.3.pipelined.144.T}{1.05}
\setval{pst.3.pipelined.18.E}{0.38}
\setval{pst.3.pipelined.18.S}{6.8}
\setval{pst.3.pipelined.18.T}{4.48}
\setval{pst.3.pipelined.2.E}{0.98}
\setval{pst.3.pipelined.2.S}{2.0}
\setval{pst.3.pipelined.2.T}{15.5}
\setval{pst.3.pipelined.288.E}{0.20}
\setval{pst.3.pipelined.288.S}{56.3}
\setval{pst.3.pipelined.288.T}{0.542}
\setval{pst.3.pipelined.2node.E}{0.20}
\setval{pst.3.pipelined.2node.S}{29.0}
\setval{pst.3.pipelined.2node.T}{1.05}
\setval{pst.3.pipelined.36.E}{0.20}
\setval{pst.3.pipelined.36.S}{7.2}
\setval{pst.3.pipelined.36.T}{4.26}
\setval{pst.3.pipelined.4.E}{0.95}
\setval{pst.3.pipelined.4.S}{3.8}
\setval{pst.3.pipelined.4.T}{8.05}
\setval{pst.3.pipelined.72.E}{0.20}
\setval{pst.3.pipelined.72.S}{14.1}
\setval{pst.3.pipelined.72.T}{2.16}
\setval{pst.3.pipelined.9.E}{0.67}
\setval{pst.3.pipelined.9.S}{6.0}
\setval{pst.3.pipelined.9.T}{5.09}
\setval{pst.3.pipelined.half.E}{0.20}
\setval{pst.3.pipelined.half.S}{7.2}
\setval{pst.3.pipelined.half.T}{4.26}
\setval{pst.3.pipelined.node.E}{0.20}
\setval{pst.3.pipelined.node.S}{14.1}
\setval{pst.3.pipelined.node.T}{2.16}
\setval{pst.3.pipelined.one.E}{1.00}
\setval{pst.3.pipelined.one.S}{1.0}
\setval{pst.3.pipelined.one.T}{30.5}
\setval{pst.3.pipelined.q.E}{0.38}
\setval{pst.3.pipelined.q.S}{6.8}
\setval{pst.3.pipelined.q.T}{4.48}
\setval{pst.3.pipelined.top.E}{0.20}
\setval{pst.3.pipelined.top.S}{56.3}
\setval{pst.3.pipelined.top.T}{0.542}
\setval{pst.3.pratio.1}{1.07}
\setval{pst.3.pratio.144}{1.38}
\setval{pst.3.pratio.18}{1.31}
\setval{pst.3.pratio.2}{1.09}
\setval{pst.3.pratio.288}{1.44}
\setval{pst.3.pratio.2node}{1.38}
\setval{pst.3.pratio.36}{1.42}
\setval{pst.3.pratio.4}{1.10}
\setval{pst.3.pratio.72}{1.44}
\setval{pst.3.pratio.9}{1.15}
\setval{pst.3.pratio.half}{1.42}
\setval{pst.3.pratio.node}{1.44}
\setval{pst.3.pratio.one}{1.07}
\setval{pst.3.pratio.q}{1.31}
\setval{pst.3.pratio.top}{1.44}
\setval{pwk.2.classic.1.E}{1.00}
\setval{pwk.2.classic.1.T}{0.558}
\setval{pwk.2.classic.144.E}{0.17}
\setval{pwk.2.classic.144.T}{3.22}
\setval{pwk.2.classic.18.E}{0.32}
\setval{pwk.2.classic.18.T}{1.77}
\setval{pwk.2.classic.2.E}{0.66}
\setval{pwk.2.classic.2.T}{0.851}
\setval{pwk.2.classic.2.ns}{4.06}
\setval{pwk.2.classic.288.E}{0.16}
\setval{pwk.2.classic.288.T}{3.41}
\setval{pwk.2.classic.2node.E}{0.17}
\setval{pwk.2.classic.2node.E2}{0.26}
\setval{pwk.2.classic.2node.T}{3.22}
\setval{pwk.2.classic.2node.cbw}{472}
\setval{pwk.2.classic.2node.cns}{14.6}
\setval{pwk.2.classic.2node.cpk}{0.77}
\setval{pwk.2.classic.2node.vec}{16.8}
\setval{pwk.2.classic.36.E}{0.18}
\setval{pwk.2.classic.36.T}{3.11}
\setval{pwk.2.classic.4.E}{0.61}
\setval{pwk.2.classic.4.T}{0.909}
\setval{pwk.2.classic.72.E}{0.18}
\setval{pwk.2.classic.72.T}{3.13}
\setval{pwk.2.classic.9.E}{0.47}
\setval{pwk.2.classic.9.T}{1.19}
\setval{pwk.2.classic.cns.hi}{14.7}
\setval{pwk.2.classic.cns.lo}{14.6}
\setval{pwk.2.classic.half.E}{0.18}
\setval{pwk.2.classic.half.T}{3.11}
\setval{pwk.2.classic.node.E}{0.18}
\setval{pwk.2.classic.node.E2}{0.27}
\setval{pwk.2.classic.node.T}{3.13}
\setval{pwk.2.classic.node.cbw}{470}
\setval{pwk.2.classic.node.cns}{14.7}
\setval{pwk.2.classic.node.cpk}{0.76}
\setval{pwk.2.classic.node.vec}{16.8}
\setval{pwk.2.classic.one.E}{1.00}
\setval{pwk.2.classic.one.T}{0.558}
\setval{pwk.2.classic.one.ns}{2.66}
\setval{pwk.2.classic.q.E}{0.32}
\setval{pwk.2.classic.q.T}{1.77}
\setval{pwk.2.classic.ratio.2node.node}{1.03}
\setval{pwk.2.classic.top.E}{0.16}
\setval{pwk.2.classic.top.E2}{0.25}
\setval{pwk.2.classic.top.T}{3.41}
\setval{pwk.2.classic.top.cbw}{471}
\setval{pwk.2.classic.top.cns}{14.7}
\setval{pwk.2.classic.top.cpk}{0.77}
\setval{pwk.2.classic.top.vec}{16.8}
\setval{pwk.2.pipelined.1.E}{1.00}
\setval{pwk.2.pipelined.1.T}{0.725}
\setval{pwk.2.pipelined.144.E}{0.17}
\setval{pwk.2.pipelined.144.T}{4.33}
\setval{pwk.2.pipelined.18.E}{0.31}
\setval{pwk.2.pipelined.18.T}{2.32}
\setval{pwk.2.pipelined.2.E}{0.91}
\setval{pwk.2.pipelined.2.T}{0.799}
\setval{pwk.2.pipelined.288.E}{0.16}
\setval{pwk.2.pipelined.288.T}{4.53}
\setval{pwk.2.pipelined.2node.E}{0.17}
\setval{pwk.2.pipelined.2node.T}{4.33}
\setval{pwk.2.pipelined.36.E}{0.17}
\setval{pwk.2.pipelined.36.T}{4.27}
\setval{pwk.2.pipelined.4.E}{0.80}
\setval{pwk.2.pipelined.4.T}{0.910}
\setval{pwk.2.pipelined.72.E}{0.17}
\setval{pwk.2.pipelined.72.T}{4.28}
\setval{pwk.2.pipelined.9.E}{0.53}
\setval{pwk.2.pipelined.9.T}{1.38}
\setval{pwk.2.pipelined.half.E}{0.17}
\setval{pwk.2.pipelined.half.T}{4.27}
\setval{pwk.2.pipelined.node.E}{0.17}
\setval{pwk.2.pipelined.node.T}{4.28}
\setval{pwk.2.pipelined.one.E}{1.00}
\setval{pwk.2.pipelined.one.T}{0.725}
\setval{pwk.2.pipelined.q.E}{0.31}
\setval{pwk.2.pipelined.q.T}{2.32}
\setval{pwk.2.pipelined.top.E}{0.16}
\setval{pwk.2.pipelined.top.T}{4.53}
\setval{pwk.2.pratio.1}{1.30}
\setval{pwk.2.pratio.144}{1.34}
\setval{pwk.2.pratio.18}{1.31}
\setval{pwk.2.pratio.2}{0.94}
\setval{pwk.2.pratio.288}{1.33}
\setval{pwk.2.pratio.2node}{1.34}
\setval{pwk.2.pratio.36}{1.37}
\setval{pwk.2.pratio.4}{1.00}
\setval{pwk.2.pratio.72}{1.37}
\setval{pwk.2.pratio.9}{1.16}
\setval{pwk.2.pratio.half}{1.37}
\setval{pwk.2.pratio.node}{1.37}
\setval{pwk.2.pratio.one}{1.30}
\setval{pwk.2.pratio.q}{1.31}
\setval{pwk.2.pratio.top}{1.33}
\setval{pwk.3.classic.1.E}{1.00}
\setval{pwk.3.classic.1.T}{0.642}
\setval{pwk.3.classic.144.E}{0.17}
\setval{pwk.3.classic.144.T}{3.77}
\setval{pwk.3.classic.18.E}{0.33}
\setval{pwk.3.classic.18.T}{1.96}
\setval{pwk.3.classic.2.E}{0.68}
\setval{pwk.3.classic.2.T}{0.948}
\setval{pwk.3.classic.2.ns}{4.54}
\setval{pwk.3.classic.288.E}{0.17}
\setval{pwk.3.classic.288.T}{3.79}
\setval{pwk.3.classic.2node.E}{0.17}
\setval{pwk.3.classic.2node.E2}{0.25}
\setval{pwk.3.classic.2node.T}{3.77}
\setval{pwk.3.classic.2node.cbw}{468}
\setval{pwk.3.classic.2node.cns}{14.8}
\setval{pwk.3.classic.2node.cpk}{0.76}
\setval{pwk.3.classic.2node.vec}{16.8}
\setval{pwk.3.classic.36.E}{0.19}
\setval{pwk.3.classic.36.T}{3.40}
\setval{pwk.3.classic.4.E}{0.65}
\setval{pwk.3.classic.4.T}{0.987}
\setval{pwk.3.classic.72.E}{0.19}
\setval{pwk.3.classic.72.T}{3.41}
\setval{pwk.3.classic.9.E}{0.51}
\setval{pwk.3.classic.9.T}{1.25}
\setval{pwk.3.classic.cns.hi}{15.0}
\setval{pwk.3.classic.cns.lo}{14.8}
\setval{pwk.3.classic.half.E}{0.19}
\setval{pwk.3.classic.half.T}{3.40}
\setval{pwk.3.classic.node.E}{0.19}
\setval{pwk.3.classic.node.E2}{0.28}
\setval{pwk.3.classic.node.T}{3.41}
\setval{pwk.3.classic.node.cbw}{467}
\setval{pwk.3.classic.node.cns}{14.8}
\setval{pwk.3.classic.node.cpk}{0.76}
\setval{pwk.3.classic.node.vec}{16.8}
\setval{pwk.3.classic.one.E}{1.00}
\setval{pwk.3.classic.one.T}{0.642}
\setval{pwk.3.classic.one.ns}{3.06}
\setval{pwk.3.classic.q.E}{0.33}
\setval{pwk.3.classic.q.T}{1.96}
\setval{pwk.3.classic.ratio.2node.node}{1.11}
\setval{pwk.3.classic.top.E}{0.17}
\setval{pwk.3.classic.top.E2}{0.25}
\setval{pwk.3.classic.top.T}{3.79}
\setval{pwk.3.classic.top.cbw}{459}
\setval{pwk.3.classic.top.cns}{15.0}
\setval{pwk.3.classic.top.cpk}{0.75}
\setval{pwk.3.classic.top.vec}{16.8}
\setval{pwk.3.pipelined.1.E}{1.00}
\setval{pwk.3.pipelined.1.T}{0.770}
\setval{pwk.3.pipelined.144.E}{0.16}
\setval{pwk.3.pipelined.144.T}{4.67}
\setval{pwk.3.pipelined.18.E}{0.30}
\setval{pwk.3.pipelined.18.T}{2.57}
\setval{pwk.3.pipelined.2.E}{0.78}
\setval{pwk.3.pipelined.2.T}{0.987}
\setval{pwk.3.pipelined.288.E}{0.15}
\setval{pwk.3.pipelined.288.T}{5.08}
\setval{pwk.3.pipelined.2node.E}{0.16}
\setval{pwk.3.pipelined.2node.T}{4.67}
\setval{pwk.3.pipelined.36.E}{0.16}
\setval{pwk.3.pipelined.36.T}{4.78}
\setval{pwk.3.pipelined.4.E}{0.71}
\setval{pwk.3.pipelined.4.T}{1.09}
\setval{pwk.3.pipelined.72.E}{0.16}
\setval{pwk.3.pipelined.72.T}{4.70}
\setval{pwk.3.pipelined.9.E}{0.53}
\setval{pwk.3.pipelined.9.T}{1.44}
\setval{pwk.3.pipelined.half.E}{0.16}
\setval{pwk.3.pipelined.half.T}{4.78}
\setval{pwk.3.pipelined.node.E}{0.16}
\setval{pwk.3.pipelined.node.T}{4.70}
\setval{pwk.3.pipelined.one.E}{1.00}
\setval{pwk.3.pipelined.one.T}{0.770}
\setval{pwk.3.pipelined.q.E}{0.30}
\setval{pwk.3.pipelined.q.T}{2.57}
\setval{pwk.3.pipelined.top.E}{0.15}
\setval{pwk.3.pipelined.top.T}{5.08}
\setval{pwk.3.pratio.1}{1.20}
\setval{pwk.3.pratio.144}{1.24}
\setval{pwk.3.pratio.18}{1.31}
\setval{pwk.3.pratio.2}{1.04}
\setval{pwk.3.pratio.288}{1.34}
\setval{pwk.3.pratio.2node}{1.24}
\setval{pwk.3.pratio.36}{1.41}
\setval{pwk.3.pratio.4}{1.10}
\setval{pwk.3.pratio.72}{1.38}
\setval{pwk.3.pratio.9}{1.15}
\setval{pwk.3.pratio.half}{1.41}
\setval{pwk.3.pratio.node}{1.38}
\setval{pwk.3.pratio.one}{1.20}
\setval{pwk.3.pratio.q}{1.31}
\setval{pwk.3.pratio.top}{1.34}
\setval{rankcounts}{1, 2, 4, 9, 18, 36, 72, 144, 256 and 288}
\setval{rec.cells}{329}
\setval{rep.cells}{247}
\setval{rep.long.max}{0.5\,\%}
\setval{rep.long.min}{5}
\setval{rep.long.n}{2}
\setval{rep.median}{1.5\,\%}
\setval{rep.nthr}{81}
\setval{rep.serial.max1s}{48\,\%}
\setval{rep.thr}{5\,\%}
\setval{repeats}{3}
\setval{scale.beyond}{, and beyond that the ranks of different nodes communicate over the network}
\setval{scale.maxcores}{288}
\setval{scale.note}{}
\setval{scale.where}{They use up to 4 whole nodes of 72 cores, 288 cores in all.}
\setval{ser.2.n1024.N}{1{,}048{,}576}
\setval{ser.2.n1024.ns}{3.57}
\setval{ser.2.n128.N}{16{,}384}
\setval{ser.2.n128.ns}{2.60}
\setval{ser.2.n2048.N}{4{,}194{,}304}
\setval{ser.2.n2048.ns}{4.09}
\setval{ser.2.n256.N}{65{,}536}
\setval{ser.2.n256.ns}{3.19}
\setval{ser.2.n4096.N}{16{,}777{,}216}
\setval{ser.2.n4096.ns}{3.49}
\setval{ser.2.n512.N}{262{,}144}
\setval{ser.2.n512.ns}{3.73}
\setval{ser.2.n64.N}{4{,}096}
\setval{ser.2.n64.ns}{6.82}
\setval{ser.2.n8192.N}{67{,}108{,}864}
\setval{ser.2.n8192.ns}{3.59}
\setval{ser.3.n101.N}{1{,}030{,}301}
\setval{ser.3.n101.ns}{3.92}
\setval{ser.3.n16.N}{4{,}096}
\setval{ser.3.n16.ns}{9.58}
\setval{ser.3.n161.N}{4{,}173{,}281}
\setval{ser.3.n161.ns}{4.18}
\setval{ser.3.n25.N}{15{,}625}
\setval{ser.3.n25.ns}{2.56}
\setval{ser.3.n256.N}{16{,}777{,}216}
\setval{ser.3.n256.ns}{3.71}
\setval{ser.3.n40.N}{64{,}000}
\setval{ser.3.n40.ns}{2.56}
\setval{ser.3.n406.N}{66{,}923{,}416}
\setval{ser.3.n406.ns}{4.24}
\setval{ser.3.n64.N}{262{,}144}
\setval{ser.3.n64.ns}{4.15}
\setval{ser.bpnnz2}{19.2}
\setval{ser.itslope2}{1.02}
\setval{ser.itslope3}{1.04}
\setval{ser.line2}{66}
\setval{ser.ns2}{3.59}
\setval{ser.ns3}{4.24}
\setval{ser.nsizes2}{8}
\setval{ser.nsizes3}{8}
\setval{ser.nsnnz2}{0.72}
\setval{ser.nsnnz3}{0.61}
\setval{ser.plane3}{1.32}
\setval{ser.ratio32}{1.18}
\setval{ser.slope2}{0.98}
\setval{ser.slope3}{0.98}
\setval{ser.tcslope2}{0.10}
\setval{ser.tcslope3}{0.09}
\setval{ser.tslope.pred2}{1.50}
\setval{ser.tslope.pred3}{1.33}
\setval{ser.tslope2}{1.61}
\setval{ser.tslope3}{1.43}
\setval{ser.ws.2.n1024}{33.6}
\setval{ser.ws.2.n128}{0.5}
\setval{ser.ws.2.n2048}{134.2}
\setval{ser.ws.2.n256}{2.1}
\setval{ser.ws.2.n4096}{536.9}
\setval{ser.ws.2.n512}{8.4}
\setval{ser.ws.2.n64}{0.1}
\setval{ser.ws.2.n8192}{2{,}147.5}
\setval{ser.ws.3.n101}{33.0}
\setval{ser.ws.3.n16}{0.1}
\setval{ser.ws.3.n161}{133.5}
\setval{ser.ws.3.n25}{0.5}
\setval{ser.ws.3.n256}{536.9}
\setval{ser.ws.3.n40}{2.0}
\setval{ser.ws.3.n406}{2{,}141.5}
\setval{ser.ws.3.n64}{8.4}
\setval{serial.Nmax}{$6.7\times10^{7}$}
\setval{serial.Nmin}{$4.1\times10^{3}$}
\setval{serial.list2}{64, 128, 256, 512, 1024, 2048, 4096 and 8192}
\setval{serial.list3}{16, 25, 40, 64, 101, 161, 256 and 406}
\setval{serialtol.list2}{64, 128, 256, 512 and 1024}
\setval{serialtol.list3}{16, 25, 40, 64 and 101}
\setval{serialtol.tol}{$1\times10^{-8}$}
\setval{slab.max}{2}
\setval{slab.thick}{1.4}
\setval{spec.cpu.cache}{67.5}
\setval{spec.cpu.cachecore}{1.88}
\setval{spec.cpu.channels}{8}
\setval{spec.cpu.memtype}{DDR5}
\setval{spec.cpu.mts}{4{,}800}
\setval{spec.cpu.peak}{307.2}
\setval{spec.gpu.bw}{696}
\setval{spec.gpu.mem}{48}
\setval{spec.gpu.name}{A40}
\setval{spec.gpu.vsnode}{1.13}
\setval{spec.node.peak}{614.4}
\setval{spread.cap}{}
\setval{spread.except}{ The one exception is a run at 256 cores on 8 nodes $\times$ 32 cores that shared its nodes with other jobs.}
\setval{spread.sent}{}
\setval{spread.tabnote}{}
\setval{st.2.classic.Enode}{0.25}
\setval{st.2.classic.Etop}{0.25}
\setval{st.2.classic.Stop}{72.2}
\setval{st.2.classic.T1}{24.1}
\setval{st.2.classic.Ttop}{0.334}
\setval{st.2.pipelined.Enode}{0.19}
\setval{st.2.pipelined.Etop}{0.20}
\setval{st.2.pipelined.Stop}{56.7}
\setval{st.2.pipelined.T1}{26.5}
\setval{st.2.pipelined.Ttop}{0.467}
\setval{st.3.classic.Enode}{0.26}
\setval{st.3.classic.Etop}{0.26}
\setval{st.3.classic.Stop}{75.3}
\setval{st.3.classic.T1}{28.4}
\setval{st.3.classic.Ttop}{0.377}
\setval{st.3.pipelined.Enode}{0.20}
\setval{st.3.pipelined.Etop}{0.20}
\setval{st.3.pipelined.Stop}{56.3}
\setval{st.3.pipelined.T1}{30.5}
\setval{st.3.pipelined.Ttop}{0.542}
\setval{st.counts}{1, 2, 4, 9, 18, 36, 72, 144 and 288}
\setval{strong.N2}{67{,}108{,}864}
\setval{strong.N3}{66{,}923{,}416}
\setval{strong.Nsci2}{$6.7\times10^{7}$}
\setval{strong.Nsci3}{$6.7\times10^{7}$}
\setval{strong.mem2}{4.3}
\setval{strong.mem3}{4.3}
\setval{strong.n2}{8{,}192}
\setval{strong.n3}{406}
\setval{strong.nnz2}{335{,}511{,}552}
\setval{strong.nnz3}{467{,}474{,}896}
\setval{topcores}{288}
\setval{topnodes}{4}
\setval{val.agree.iters}{25}
\setval{val.agree.margin}{100}
\setval{val.agree.max}{$8.4\times10^{-13}$}
\setval{val.agree.max2}{$4.9\times10^{-14}$}
\setval{val.agree.max3}{$8.4\times10^{-13}$}
\setval{val.agree.tol}{$1\times10^{-10}$}
\setval{val.bug.max}{0.76}
\setval{val.bug.min}{0.46}
\setval{val.bug.n}{2}
\setval{val.cl.sha}{9da57b11ddd2, a0006417c428, ace45bed3220}
\setval{val.cluster}{vanda}
\setval{val.cond.2}{1{,}659}
\setval{val.cond.3}{414}
\setval{val.floor.classic.iters}{258}
\setval{val.floor.classic.rec}{$9.9\times10^{-15}$}
\setval{val.floor.classic.true}{$5.3\times10^{-13}$}
\setval{val.floor.classic.verdict}{reached it in 258 iterations, with a true residual of $5.3\times10^{-13}$}
\setval{val.floor.maxit}{3{,}000}
\setval{val.floor.n}{63}
\setval{val.floor.pipe.best}{$1.1\times10^{-10}$}
\setval{val.floor.pipe.sent}{ In V5 the same case had reached a true residual of $1.1\times10^{-10}$ after 215 iterations, so in the longer run the true residual rose while the tracked one kept falling.}
\setval{val.floor.pipelined.iters}{3{,}000}
\setval{val.floor.pipelined.rec}{$2.9\times10^{-12}$}
\setval{val.floor.pipelined.true}{$1.5\times10^{-9}$}
\setval{val.floor.pipelined.verdict}{did not reach it in 3{,}000 iterations: its tracked residual had fallen only to $2.9\times10^{-12}$, and its true residual was $1.5\times10^{-9}$ at the end}
\setval{val.floor.ranks}{4}
\setval{val.floor.tol}{$1\times10^{-14}$}
\setval{val.ncases}{5}
\setval{val.order.V1}{1.998 and 2.000}
\setval{val.order.V2}{1.989 and 2.000}
\setval{val.order.hi}{2.1}
\setval{val.order.lo}{1.9}
\setval{val.ordertol}{$1\times10^{-12}$}
\setval{val.pipe.iters}{the same number of iterations as classic CG}
\setval{val.pipe.maxdiff}{$1.2\times10^{-12}$}
\setval{val.pipe.minbound}{$7.7\times10^{-8}$}
\setval{val.pipe.truetol}{$1\times10^{-9}$}
\setval{val.pipetol}{$1\times10^{-10}$}
\setval{val.programs}{The CPU solver is \code{a0006417c428}, validated with V1--V5 in run 20261112T063938Z. The graph solver is \code{9da57b11ddd2}, validated with G1 and G2 in run 20261112T072321Z. The GPU version is \code{ace45bed3220}, validated with V1--V5 in run 20261112T225342Z on the GPU node that then timed it.}
\setval{val.runs}{27}
\setval{val.sha}{a0006417c428}
\setval{val.stamp}{20261112T063938Z}
\setval{val.timed.sent}{Every program behind a timed result passed its validation tests on Vanda before its first timed run (Section~\ref{sec:submit} lists the programs).}
\setval{val.where}{a Vanda compute node}
\setval{var.2.node.best}{2D decomposition with classic CG and blocking halo exchange}
\setval{var.2.node.classic.spread}{0.2\,\%}
\setval{var.2.node.spread}{1.39}
\setval{var.2.top.best}{1D decomposition with classic CG and blocking halo exchange}
\setval{var.2.top.spread}{1.52}
\setval{var.3.node.best}{2D decomposition with classic CG and blocking halo exchange}
\setval{var.3.node.blocks}{1.027}
\setval{var.3.node.d1.comp}{14.94}
\setval{var.3.node.d1.halo}{3.23}
\setval{var.3.node.d1.red}{1.52}
\setval{var.3.node.d2.comp}{13.56}
\setval{var.3.node.d2.halo}{0.69}
\setval{var.3.node.d2.red}{0.32}
\setval{var.3.node.d3.comp}{13.50}
\setval{var.3.node.d3.halo}{0.67}
\setval{var.3.node.d3.red}{0.79}
\setval{var.3.node.slab}{1.35}
\setval{var.3.node.spread}{1.74}
\setval{var.3.top.best}{2D decomposition with classic CG and non-blocking halo exchange}
\setval{var.3.top.d1.comp}{3.31}
\setval{var.3.top.d1.halo}{3.19}
\setval{var.3.top.d1.red}{2.48}
\setval{var.3.top.d2.comp}{2.63}
\setval{var.3.top.d2.halo}{0.50}
\setval{var.3.top.d2.red}{0.67}
\setval{var.3.top.d3.comp}{2.73}
\setval{var.3.top.d3.halo}{0.47}
\setval{var.3.top.d3.red}{0.56}
\setval{var.3.top.slab}{2.36}
\setval{var.3.top.spread}{2.93}
\setval{var.node.ovl.gain}{0.3\,\%}
\setval{var.node.ovl.loss}{5.7\,\%}
\setval{var.top.ovl.gain}{4.1\,\%}
\setval{var.top.ovl.loss}{11.0\,\%}
\setval{vec.classic}{5}
\setval{vec.extra}{3}
\setval{vec.pipelined}{8}
\setval{weak.n2}{1{,}448}
\setval{weak.n3}{128}
\setval{weak.per2}{2{,}096{,}704}
\setval{weak.per3}{2{,}097{,}152}
\setval{weak.top.N2}{$6.0\times10^{8}$}
\setval{weak.top.N3}{$6.0\times10^{8}$}
\setval{weak.top.n2}{24{,}573}
\setval{weak.top.n3}{845}
\setval{weak.ws2}{67}
\setval{weak.ws3}{67}
\setval{wk.2.classic.Enode}{0.18}
\setval{wk.2.classic.Etop}{0.16}
\setval{wk.2.classic.T1}{0.558}
\setval{wk.2.classic.Ttop}{3.41}
\setval{wk.2.pipelined.Enode}{0.17}
\setval{wk.2.pipelined.Etop}{0.16}
\setval{wk.2.pipelined.T1}{0.725}
\setval{wk.2.pipelined.Ttop}{4.53}
\setval{wk.3.classic.Enode}{0.19}
\setval{wk.3.classic.Etop}{0.17}
\setval{wk.3.classic.T1}{0.642}
\setval{wk.3.classic.Ttop}{3.79}
\setval{wk.3.pipelined.Enode}{0.16}
\setval{wk.3.pipelined.Etop}{0.15}
\setval{wk.3.pipelined.T1}{0.770}
\setval{wk.3.pipelined.Ttop}{5.08}
\setval{wk.top.n2}{24{,}573}
\setval{wk.top.n3}{845}

% generated by scripts/report_macros.py; do not edit
\setslot{fig_serial}{1}{}
\setslot{fig_serial_tol}{1}{}
\setslot{fig_strong}{1}{}
\setslot{fig_weak}{1}{}
\setslot{fig_efficiency}{1}{}
\setslot{fig_comm}{1}{}
\setslot{fig_layout}{1}{}
\setslot{fig_graph_serial}{1}{}
\setslot{fig_graph_strong}{1}{}
\setslot{fig_graph_comm}{1}{}
\setslot{fig_graph_variants}{1}{}
\setslot{fig_graph_sweep}{1}{}

\lstset{basicstyle=\ttfamily\scriptsize,columns=fullflexible,keepspaces=true,breaklines=true,
  breakatwhitespace=false,numbers=left,numberstyle=\tiny\color{gray},numbersep=6pt,xleftmargin=1.6em,
  showstringspaces=false,tabsize=4,upquote=true}
% generated by scripts/report_macros.py: the list of appendix files (printed from ASCII copies)
\newcommand{\appendixfiles}{src/par.hpp,src/grid.hpp,src/grid.cpp,src/kernels.hpp,src/kernels.cpp,src/halo.hpp,src/halo.cpp,src/cg.hpp,src/cg.cpp,src/record.hpp,src/record.cpp,src/main.cpp,src/cg_impl.hpp,src/graph/csrfile.hpp,src/graph/dcsr.hpp,src/graph/dcsr.cpp,src/graph/label.hpp,src/graph/label.cpp,src/graph/graph_main.cpp,src/graph/convert.cpp,CMakeLists.txt,scripts/build.sh,scripts/fetch_snap.sh,scripts/convert_snap.sh,cluster/vanda/env.sh,checks/validate.py,checks/validate_graph.py,README.md}


\begin{document}

\begin{center}
{\LARGE\bfseries\sffamily Conjugate gradient for the Poisson equation\\[2pt] in two and three dimensions with MPI and OpenMP}\\[8pt]
DSA5208 Project 3\\[3pt]
Kushagra Bhatnagar (A0351591A) \quad$\cdot$\quad Pratyush Sinha (A0352482A)
\end{center}

\textbf{Summary.} We solve the Poisson equation $-\Delta u = f$ on the unit square and the unit cube with the
\term{conjugate gradient} (CG) method, an iterative solver for symmetric positive definite linear systems, without
ever storing the matrix. It is parallelised with MPI across processes and with OpenMP \term{threads} (workers that
share one process's memory) inside each process, and the same loops can also run on a GPU through OpenMP GPU offload.
The parallel solver offers three switchable choices: how the grid is split (along one, two or three directions),
whether boundary values are exchanged between processes before or during the computation, and whether CG waits for two
reductions per iteration or overlaps a single one with the work. A small set of \val{val.ncases} validation tests
shows that the error falls as $h^2$ and that every parallel configuration reproduces the serial result to rounding,
and the tests fail on two deliberately broken versions of the boundary exchange. On \val{cl.name}, with the problem fixed
and the cores increased to \val{meas.top}\val{meas.where}, classic CG keeps a \term{parallel efficiency} (how close the run
comes to being $P$ times as fast on $P$ cores) of \val{st.3.classic.Etop} in 3D and \val{st.2.classic.Etop} in 2D. With the work per core fixed,
it keeps \val{wk.3.classic.Etop} (3D) and \val{wk.2.classic.Etop} (2D) of the one-core speed.\val{scale.note} At \val{meas.top}
cores, exchanging boundary values takes \val{cm.3.halo} and reductions \val{cm.3.red} of the 3D run time. One GPU
runs the 3D problem \val{gpu.3.vsnode} times as fast as the best layout on one \val{lay.cpn}-core node of
\val{cl.layout}.

Section~\ref{sec:algorithm} states the method and the synthetic test problem, Section~\ref{sec:design} the parallel
design and its cost model, and Section~\ref{sec:validation} the validation tests. Section~\ref{sec:env} describes the
nodes and how each time is measured, and the results follow: the serial cost in Section~\ref{sec:serial}, MPI scaling
in Section~\ref{sec:mpi}, and MPI with OpenMP threads and the GPU in Section~\ref{sec:mpix}.
Section~\ref{sec:graph} applies the same solver to real graphs, and Section~\ref{sec:discussion} draws the results
together.

% ======================================================================================================
\section{The algorithm and the synthetic test problem}\label{sec:algorithm}
\req{E1}\req{E12}\req{E11}\req{E13}
\subsection{The Poisson problem and its discretisation}
We solve $-\Delta u = f$ on $\Omega=(0,1)^d$ for $d=2$ and $d=3$ with $u=0$ on the boundary, and the same program
handles both dimensions, as the project asks for both. A uniform grid with $n$ interior points per direction and
spacing $h=1/(n+1)$ covers the domain, and the second-order finite-difference stencil replaces the Laplacian, the
5-point stencil in 2D and the 7-point stencil in 3D:
\begin{equation*}
\frac{1}{h^2}\Big(2d\,u_{\mathbf{i}} - \sum_{a=1}^{d}\big(u_{\mathbf{i}+\mathbf{e}_a} + u_{\mathbf{i}-\mathbf{e}_a}\big)\Big)
= f_{\mathbf{i}},
\end{equation*}
where $\mathbf{e}_a$ is the unit step in direction $a$ and a neighbour outside the domain contributes its boundary
value, zero. Collecting the equations gives a symmetric positive definite system $A\mathbf{u}=\mathbf{b}$ with $N=n^d$
unknowns and $n^d + 2d(n-1)n^{d-1}$ nonzero entries, slightly fewer than $5N$ in 2D and $7N$ in 3D because the rows
next to the boundary lose neighbours. The solution of this system differs from the exact solution by a discretisation
error that falls as $h^2$. We multiply both sides by $h^2$, which leaves the solution unchanged and makes the entries
of $A$ small integers.

The solver never stores $A$. A product with $A$ is one pass over the grid that reads each point and its $2d$
neighbours, and a solver that works only through such products is called \term{matrix-free}. We chose this because
the stencil describes $A$ completely, so storing it would add memory traffic and nothing else. One iteration of
classic CG performs \val{cg.flops2} floating-point operations per grid point in 2D and \val{cg.flops3} in 3D, counted in
the loops of the program, and moves about \val{bytes.classic} bytes per point between memory and the processor
(Section~\ref{sec:model}). Its arithmetic intensity, the operations per byte moved, is therefore only \val{cg.ai2} in
2D and \val{cg.ai3} in 3D, and at so few operations per byte the time of an iteration is set by how fast memory
delivers the vectors rather than by arithmetic. The computation is \term{memory-bound}, which the serial measurements of
Section~\ref{sec:serial} confirm.

\subsection{The conjugate gradient method}
CG solves a symmetric positive definite system by building, one iteration at a time, search directions that are
conjugate with respect to $A$, so that each iteration minimises the error in the $A$-norm over all directions found
so far. Algorithm~\ref{alg:cg} is the textbook form of Hestenes and Stiefel~\cite{hs}, started from
$\mathbf{x}_0=0$: each iteration applies $A$ once, computes two inner products and updates three vectors.

\begin{figure}[H]\centering\small
\fbox{\begin{minipage}{0.86\linewidth}
\algotitle{alg:cg}{classic CG}\\[2pt]
$\mathbf{r}=\mathbf{b}$, \ $\mathbf{p}=\mathbf{r}$, \ $\gamma=(\mathbf{r},\mathbf{r})$ \hfill\emph{one reduction}\\
repeat:\\
\hspace*{1.2em}$\mathbf{q}=A\mathbf{p}$ \hfill\emph{halo exchange, stencil}\\
\hspace*{1.2em}$\alpha=\gamma/(\mathbf{p},\mathbf{q})$ \hfill\emph{reduction (\code{MPI\_Allreduce})}\\
\hspace*{1.2em}$\mathbf{x}\leftarrow\mathbf{x}+\alpha\mathbf{p}$, \ $\mathbf{r}\leftarrow\mathbf{r}-\alpha\mathbf{q}$,
\ $\gamma'=(\mathbf{r},\mathbf{r})$ \hfill\emph{reduction (\code{MPI\_Allreduce})}\\
\hspace*{1.2em}stop if $\sqrt{\gamma'}\le\text{tol}\,\|\mathbf{b}\|$, \ otherwise
$\mathbf{p}\leftarrow\mathbf{r}+(\gamma'/\gamma)\,\mathbf{p}$, \ $\gamma\leftarrow\gamma'$
\end{minipage}}
\end{figure}

The number of iterations is governed by the condition number of $A$, here $\kappa=\cot^2(\pi h/2)\approx
\big(2(n+1)/\pi\big)^2$. The error shrinks by a fixed factor about every $\sqrt{\kappa}$ iterations, so reaching a fixed
tolerance takes a number of iterations proportional to $n$, and because one iteration costs work proportional to $N$,
the serial time to a fixed tolerance grows like $N^{1+1/d}$, that is like $N^{3/2}$ in 2D and $N^{4/3}$ in 3D.
Parallelisation changes the cost of one iteration but not the number of iterations, which is a property of the method.
We therefore time a fixed number of iterations in the scaling studies (Section~\ref{sec:timing}) and measure the growth
of the iteration count separately, in serial (Section~\ref{sec:serial}).

\subsection{The synthetic test problem}\label{sec:problem}
The Poisson part uses no dataset. Its synthetic test problem starts from a chosen exact solution $u$, from which we
compute the $f$ that produces it, so that the error of every run can be measured. This construction is called a
\term{manufactured solution}. We use
\begin{equation*}
u(\mathbf{x})=\prod_{a=1}^{d} g(x_a),\quad g(t)=t(1-t)e^{t},\qquad
f(\mathbf{x})=\sum_{a=1}^{d} x_a(x_a+3)e^{x_a}\prod_{b\ne a} g(x_b),
\end{equation*}
which vanishes on the boundary. Each process evaluates $h^2 f$ only at the grid points it owns, so a problem of any size
is generated in place and in parallel at the start of a run, and the runs in which the grid grows with the number of
processes cost nothing extra to set up.

The obvious choice $u=\prod_a\sin(\pi x_a)$ is unsuitable because it is an eigenvector of the discrete operator. The
right-hand side is then a multiple of an eigenvector, so CG finds the exact discrete solution in a single iteration,
which says nothing about convergence and leaves no work to time. The program keeps it behind an option only to show
this.

\subsection{Why the method is not embarrassingly parallel}
Every CG iteration depends on the whole grid in two ways. The stencil at a point needs the values of its neighbours, so
a process that owns part of the grid needs, in every iteration, the boundary values owned by the processes next to it.
Each inner product needs a contribution from every point, so every process takes part in a reduction
(\code{MPI\_Allreduce}) whose result all of them need before the next step can start, and Algorithm~\ref{alg:cg} has
two such reductions per iteration. Neither dependence can be removed without changing the method, and both recur in
every iteration. Section~\ref{sec:model} weighs their cost against the memory traffic of an iteration and says where
each of them limits the scaling.

% ======================================================================================================
\section{Parallel design of the solver}\label{sec:design}
\req{E8}
\subsection{Domain decomposition}
We use MPI with one rank for each box of the grid, which is domain decomposition in the sense of Lecture~12: each rank
updates its own box and exchanges data only where boxes touch. \code{MPI\_Dims\_create} and \code{MPI\_Cart\_create}
arrange the ranks in a Cartesian process grid, which gives each rank its coordinates and its neighbour in each
direction. The points of each direction are dealt out in contiguous blocks, and when $n$ is not divisible by the number
of ranks $p$ in that direction, the first $n \bmod p$ ranks get one extra point. A run chooses how many directions to
split with the option \code{-{}-decomp}: splitting one direction gives \term{slabs} (a 1D decomposition), splitting two
gives \term{columns} (2D) and splitting all three gives \term{blocks} (3D).

Each rank stores its box with one extra layer of points on every side, and this layer, which holds copies of the
neighbours' boundary values, is called the \term{ghost layer}. Before every stencil application each rank sends its
outermost layer of owned values, its halo, to the neighbouring rank, which stores it in its ghost layer, and this copy
is the halo exchange. Ghost points on the outer boundary of the domain are never written and stay zero, which is how
the boundary condition enters. Because the stencil has no diagonal neighbours, only faces are exchanged, never edges or
corners, so a rank sends at most two messages per split direction.

The decomposition decides how much each rank sends. A 3D slab exchanges two full planes however many ranks share the
grid, while the faces of a block shrink as the number of ranks grows. Table~\ref{tab:geometry} works this out for the
strong-scaling grids. At \val{topcores} ranks a 3D slab is only \val{slab.thick} planes thick on average and at most
\val{slab.max}, so its halo is \val{geo.3d.slab.ratio} of the points it owns, whereas the halo of a block is
\val{geo.3d.cube.ratio} of its points, at the price of six messages instead of two. We kept all three decompositions so
that the measured times can show where fewer, larger messages beat more, smaller ones.

\begin{table}[H]\centering\footnotesize
%% generated by scripts/report_macros.py; do not edit
\begin{tabular}{llrlrrrr}
\toprule
grid & split & ranks & process grid & points per rank & messages & halo values & halo / points \\
\midrule
2D & 1D & 288 & $288\times1$ & 237{,}568 & 2 & 16{,}384 & 6.90\,\% \\
2D & 2D & 72 & $12\times6$ & 932{,}978 & 4 & 4{,}098 & 0.44\,\% \\
2D & 2D & 288 & $24\times12$ & 233{,}586 & 4 & 2{,}050 & 0.88\,\% \\
3D & 1D & 288 & $288\times1\times1$ & 329{,}672 & 2 & 329{,}672 & 100.00\,\% \\
3D & 3D & 72 & $6\times4\times3$ & 936{,}360 & 6 & 59{,}772 & 6.38\,\% \\
3D & 3D & 288 & $8\times6\times6$ & 235{,}824 & 6 & 23{,}120 & 9.80\,\% \\
\bottomrule
\end{tabular}

\caption{Halo per rank and exchange for the strong-scaling grids (2D $n=\val{strong.n2}$, 3D $n=\val{strong.n3}$),
computed from the geometry for an interior rank with the largest box.}
\label{tab:geometry}
\end{table}

\subsection{Blocking and non-blocking halo exchange}
The blocking halo exchange calls \code{MPI\_Sendrecv} twice per split direction, first towards the higher neighbour and
then towards the lower one, and applies the stencil only when all ghost layers are filled, so the rank does nothing
while it waits. The \term{non-blocking halo exchange} instead starts all receives and sends with \code{MPI\_Irecv} and
\code{MPI\_Isend}, applies the stencil to the interior points, whose neighbours are all local, waits with
\code{MPI\_Waitall}, and then applies the stencil to the outer shell of points that need ghost values. Only the points
next to a face shared with another rank are computed after the wait, so a box on the edge of the domain has a shell on
fewer sides. The aim is to hide the
message time behind computation, but whether it is hidden depends on the MPI library, because many implementations move
a large message forward only while the program is inside an MPI call, which here happens only at the final wait. We
therefore measure the overlap rather than assume it (Section~\ref{sec:mpi}). Both exchanges use two-sided messages, and
we did not implement the one-sided alternative with MPI windows.

\subsection{Classic CG and a pipelined variant}
Classic CG waits for two reductions per iteration, and each lasts at least the latency of one message times the depth
of the reduction tree, about $\log_2 P$ steps for $P$ ranks, however little work each rank has.

Ghysels and Vanroose~\cite{gv} rearrange the recurrences of CG so that each iteration needs a single reduction of two
numbers, which is started with the non-blocking \code{MPI\_Iallreduce} and runs while the next matrix--vector product is
computed (Algorithm~\ref{alg:pcg}), and we call this form \term{pipelined CG}. In exact arithmetic it produces the same
iterates as classic CG. It keeps \val{vec.pipelined} vectors instead of \val{vec.classic} and updates most of them in
one loop, and counting the vectors that each loop reads and writes (Section~\ref{sec:model}) gives about
\val{bytes.pipelined} bytes per grid point and iteration against \val{bytes.classic} for classic CG. On a memory-bound
problem an iteration of pipelined CG therefore takes longer unless there is reduction latency to hide. Its residual is also updated through a longer chain of recurrences, so rounding errors accumulate
further: the residual that the method tracks, which we call the \term{recursive residual}, drifts away from the
\term{true residual} $\mathbf{b}-A\mathbf{x}$, and this limits the accuracy that pipelined CG can reach
(Section~\ref{sec:validation}). We kept both variants because only measurement on the target machine can decide whether
the shorter wait outweighs the extra memory traffic.

\begin{figure}[H]\centering\small
\fbox{\begin{minipage}{0.86\linewidth}
\algotitle{alg:pcg}{pipelined CG, without preconditioner}\\[2pt]
$\mathbf{r}=\mathbf{b}$, \ $\mathbf{w}=A\mathbf{r}$\\
repeat:\\
\hspace*{1.2em}start $\gamma=(\mathbf{r},\mathbf{r})$, $\delta=(\mathbf{w},\mathbf{r})$ \hfill\emph{one non-blocking
reduction (\code{MPI\_Iallreduce})}\\
\hspace*{1.2em}$\mathbf{m}=A\mathbf{w}$ \hfill\emph{halo exchange, stencil, while the reduction runs}\\
\hspace*{1.2em}wait for $\gamma,\delta$, and stop if $\sqrt{\gamma}\le\text{tol}\,\|\mathbf{b}\|$\\
\hspace*{1.2em}$\beta=\gamma/\gamma_{\text{old}}$, $\alpha=\gamma/(\delta-\beta\gamma/\alpha_{\text{old}})$ \
(first iteration: $\beta=0$, $\alpha=\gamma/\delta$)\\
\hspace*{1.2em}$\mathbf{z}\leftarrow\mathbf{m}+\beta\mathbf{z}$, $\mathbf{s}\leftarrow\mathbf{w}+\beta\mathbf{s}$,
$\mathbf{p}\leftarrow\mathbf{r}+\beta\mathbf{p}$, $\mathbf{x}\leftarrow\mathbf{x}+\alpha\mathbf{p}$,
$\mathbf{r}\leftarrow\mathbf{r}-\alpha\mathbf{s}$, $\mathbf{w}\leftarrow\mathbf{w}-\alpha\mathbf{z}$
\hfill\emph{one pass}
\end{minipage}}
\end{figure}

\subsection{OpenMP threads and OpenMP GPU offload}
Inside each rank, OpenMP threads split every loop over the grid, which covers the stencil, the vector updates and the
inner products, with a static schedule, so that each thread always touches the same rows. Linux places a memory page on
the socket of the thread that first writes it, a policy known as \term{first touch}, so we write every vector first with
the threads that later use it, which places each page next to the cores that read it. Only the master thread calls MPI,
and only outside the threaded loops, which is the thread level that MPI calls \code{MPI\_THREAD\_FUNNELED}, the weakest
level that allows threads at all. Fewer ranks with more threads mean fewer and larger halo messages and fewer ranks in
each reduction, but they also pay for starting and joining the threads in every loop, and Section~\ref{sec:mpix}
compares these layouts on a \val{cpn}-core node.

For the GPU we kept the same loops and added OpenMP GPU offload directives (\code{target teams distribute parallel
for}), switched on at build time. The vectors are copied to the GPU once and stay there, and in every halo exchange the
halos are packed on the GPU, copied to the host, exchanged with MPI and copied back. One rank drives one GPU. We chose
OpenMP offload over CUDA so that one source builds both versions, which then run the same loops, and so that the
program could be developed and validated on a machine without a GPU. The halos pass through the host because our MPI
calls take host buffers, and a CUDA version would copy them in the same way unless its MPI library could send directly
from GPU memory.

\subsection{The cost model}\label{sec:model}
Lecture~12 splits the run time into three parts: the floating-point operations times the time per operation, the
volume of communication divided by the bandwidth, and the number of communication calls times the latency. For our
memory-bound solver the first part is the memory traffic, which we count per grid point and iteration from the loops of
the program. A loop makes one pass through memory, which we call a stream, for each vector that it reads or writes,
and two for a vector that it writes without reading, because the cache loads a line from memory before it writes it.
Classic CG then makes \val{cg.st.classic} streams per point, \val{cg.st.classic.stencil} in the stencil,
\val{cg.st.classic.xr} in the update of $\mathbf{x}$ and $\mathbf{r}$ and \val{cg.st.classic.p} in the update of
$\mathbf{p}$, which at eight bytes each are the \val{bytes.classic} bytes used throughout. Pipelined CG makes
\val{cg.st.pipelined.stencil} in the stencil and \val{cg.st.pipelined.upd} in its single update loop, \val{bytes.pipelined}
bytes in all, a ratio of \val{cg.byteratio}.

These bytes set a floor under the time of an iteration. With $\beta$ bytes per point, $N$ points, $P$ ranks on $s$
sockets, a memory bandwidth $B$ per socket and a bandwidth $b$ that one core can draw on its own,
\begin{equation*}
T_{\text{mem}}\ \ge\ \max\Big(\frac{\beta N}{s\,B},\ \frac{\beta N}{P\,b}\Big).
\end{equation*}
The second part of the cost is the halo of Table~\ref{tab:geometry} at eight bytes per value, and the third is the halo
messages, two per split direction, plus the reductions, two per iteration for classic CG and one for pipelined CG, each
of about $\log_2 P$ latency steps.

With the problem fixed and $P$ growing, the model predicts the following. While the cores of a socket leave some of its
bandwidth unused, the memory term falls like $1/P$, but once they draw all of it, only another socket or another node
makes the term smaller, so inside a socket the time stops falling well before its cores run out, and inside a node only
the second socket helps. Across nodes the memory
term falls again with the number of sockets, the halo term falls like the area of a face, $P^{-1/2}$ for 2D blocks and
$P^{-2/3}$ for 3D blocks, and the reduction term grows like $\log_2 P$. Our thesis follows from this: inside a node,
memory bandwidth limits the scaling, and across nodes the reductions take a growing share of each iteration, until
beyond some number of nodes they and not the memory traffic set the limit. Up to \val{node.cpn2} cores, two nodes, the
measurements of Section~\ref{sec:mpi} are in the first of these regimes, and for the graphs of Section~\ref{sec:graph},
whose ghost exchange is not a thin surface, communication already takes most of the time on two nodes. The halo of the
grid stays small because it is the surface of a box (Table~\ref{tab:geometry}), not because the product is
matrix-free: a stored sparse matrix distributed in the same way would exchange the same halo.

The program times three parts on every rank with \code{MPI\_Wtime}. The computation (the stencil and the vector loops)
tests the memory term. The halo exchange contains both the volume term and the latency of the halo messages, and the
reductions contain their latency and also any waiting for slower ranks (Section~\ref{sec:timing}).

% ======================================================================================================
\section{Validation tests}\label{sec:validation}
\req{E9}
We use \val{val.ncases} validation tests (Table~\ref{tab:cases}), because the program can be wrong in three different
ways and each test checks one of them: the discretisation (V1, V2), the bookkeeping of the parallel runs (V3, V4) and
the second CG variant (V5). A program is used for timing only if it passes all of them, and the program used for every
timed run on \val{cl.used} passed them on that cluster before any timed run started. The results below were obtained
on \val{val.where}, the kind of node that Section~\ref{sec:env} describes.

\begin{table}[H]\centering\small
%% generated by scripts/report_macros.py; do not edit
\begin{tabular}{lp{0.82\linewidth}}
\toprule
test & what it checks \\
\midrule
V1 & 2D, one process: the solver converges, and its error against the exact solution falls as $h^2$ \\
V2 & 3D, one process: as V1 \\
V3 & 2D, several processes (1D and 2D decompositions, uneven blocks, blocking and non-blocking halo exchange, two threads per rank): the iterate equals the one-process iterate to rounding \\
V4 & 3D, several processes: as V3, with 1D, 2D and 3D decompositions \\
V5 & pipelined CG on one and four processes reaches the same converged answer as classic CG, 2D and 3D \\
\bottomrule
\end{tabular}

\caption{The validation tests and what each checks.}
\label{tab:cases}
\end{table}

\emph{Error against the exact solution (V1, V2).} Halving $h$ should divide the discretisation error by four. We solve
to a relative residual of \val{val.ordertol}, far below the discretisation error, so that the solver's own error does
not blur the comparison, and we require the \term{observed order}, the base-two logarithm of the ratio of the errors on
two successive grids, to lie between \val{val.order.lo} and \val{val.order.hi}. The observed orders are
\val{val.order.V1} in 2D and \val{val.order.V2} in 3D, against the theoretical order of two (Table~\ref{tab:valorder}),
and the iteration counts double with each halving of $h$, as the growth of $\sqrt{\kappa}$ with $n$ predicts
(Section~\ref{sec:algorithm}).

\begin{table}[H]\centering\small
%% generated by scripts/report_macros.py; do not edit
\begin{tabular}{lrrrrr}
\toprule
grid & $n$ & $h$ & iterations & max error & observed order \\
\midrule
2D & 31 & 1/32 & 118 & $9.07\times10^{-5}$ & -- \\
2D & 63 & 1/64 & 239 & $2.27\times10^{-5}$ & 1.998 \\
2D & 127 & 1/128 & 481 & $5.68\times10^{-6}$ & 2.000 \\
3D & 15 & 1/16 & 73 & $1.56\times10^{-4}$ & -- \\
3D & 31 & 1/32 & 150 & $3.93\times10^{-5}$ & 1.989 \\
3D & 63 & 1/64 & 309 & $9.82\times10^{-6}$ & 2.000 \\
\bottomrule
\end{tabular}

\caption{V1 and V2: largest error against the exact solution, iterations and observed order, one process.}
\label{tab:valorder}
\end{table}

\emph{Parallel runs equal the serial run (V3, V4).} A parallel run should differ from the serial run only through
rounding, because once the halos are right the stencil and the vector updates perform the same operations on the same
numbers in any configuration, and only the order in which the partial sums of the inner products are added changes. We
compare the iterate after \val{val.agree.iters} iterations, before convergence, because a wrong halo could hide behind a
converged answer, and we set the tolerance on the relative difference to \val{val.agree.tol}. The largest difference
observed is \val{val.agree.max} (Table~\ref{tab:valagree}), more than \val{val.agree.margin} times smaller than the
tolerance and at the level of rounding. The configurations cover all three decompositions, uneven blocks (grid sizes
that do not divide evenly among the ranks), both halo exchanges and ranks with two threads.

\begin{table}[H]\centering\small
%% generated by scripts/report_macros.py; do not edit
\begin{tabular}{lrrrlll}
\toprule
grid & $n$ & ranks & threads & decomposition & halo exchange & difference \\
\midrule
2D & 50 & 2 & 1 & 1D & blocking & $4.6\times10^{-15}$ \\
2D & 50 & 3 & 1 & 1D & non-blocking & $5.9\times10^{-15}$ \\
2D & 50 & 4 & 1 & 2D & blocking & $5.7\times10^{-15}$ \\
2D & 50 & 6 & 1 & 2D & non-blocking & $6.1\times10^{-15}$ \\
2D & 50 & 2 & 2 & 2D & non-blocking & $4.9\times10^{-14}$ \\
3D & 21 & 4 & 1 & 1D & blocking & $1.3\times10^{-13}$ \\
3D & 21 & 4 & 1 & 2D & non-blocking & $5.9\times10^{-14}$ \\
3D & 21 & 6 & 1 & 3D & non-blocking & $2.1\times10^{-13}$ \\
3D & 21 & 8 & 1 & 3D & blocking & $5.5\times10^{-13}$ \\
3D & 21 & 8 & 1 & 3D & non-blocking & $8.4\times10^{-13}$ \\
3D & 21 & 2 & 2 & 3D & blocking & $2.2\times10^{-13}$ \\
\bottomrule
\end{tabular}

\caption{V3 and V4: relative difference between a parallel run and the one-process run after the same number of
iterations.}\label{tab:valagree}
\end{table}

\emph{Pipelined CG (V5).} Pipelined and classic CG follow different recurrences, so their iterates agree only in exact
arithmetic and can be compared only once both have converged. For any $\mathbf{x}$ the relative error is at most
$\kappa$ times the relative true residual $\rho=\|\mathbf{b}-A\mathbf{x}\|/\|\mathbf{b}\|$, and applying this to both
answers bounds their difference by $\|\mathbf{x}_p-\mathbf{x}_s\|/\|\mathbf{x}_s\|\le\kappa(\rho_p+\rho_s)/(1-\kappa
\rho_s)$, where the subscript $p$ marks the pipelined and $s$ the standard, classic variant. Both variants stop when the
residual that they track falls below \val{val.pipetol} of $\|\mathbf{b}\|$, and the true residual of the pipelined
answer must also be at most \val{val.pipe.truetol}, a looser limit because its tracked residual can drift
(Section~\ref{sec:design}). Pipelined CG took \val{val.pipe.iters}, and the largest difference, \val{val.pipe.maxdiff}, lies far inside
the smallest bound, \val{val.pipe.minbound} (Table~\ref{tab:valpipe}).

\begin{table}[H]\centering\footnotesize
%% generated by scripts/report_macros.py; do not edit
\begin{tabular}{lrrlrrrrrr}
\toprule
grid & $n$ & ranks & halo exchange & it.\ classic & it.\ pipel. & $\rho_s$ & $\rho_p$ & difference & bound \\
\midrule
2D & 63 & 1 & blocking & 215 & 215 & $8.5\times10^{-11}$ & $1.2\times10^{-10}$ & $1.2\times10^{-12}$ & $3.3\times10^{-7}$ \\
2D & 63 & 4 & non-blocking & 215 & 215 & $8.5\times10^{-11}$ & $1.1\times10^{-10}$ & $9.2\times10^{-13}$ & $3.2\times10^{-7}$ \\
3D & 31 & 1 & blocking & 132 & 132 & $9.3\times10^{-11}$ & $9.3\times10^{-11}$ & $3.6\times10^{-14}$ & $7.7\times10^{-8}$ \\
3D & 31 & 4 & non-blocking & 132 & 132 & $9.3\times10^{-11}$ & $9.3\times10^{-11}$ & $5.7\times10^{-14}$ & $7.7\times10^{-8}$ \\
\bottomrule
\end{tabular}

\caption{V5: pipelined CG against classic CG, both converged. $\rho$ is the relative true residual and the bound is
$\kappa(\rho_p+\rho_s)/(1-\kappa\rho_s)$.}\label{tab:valpipe}
\end{table}

\emph{The tests detect errors in the halo exchange.} A validation test is useful only if it fails on a wrong program,
so we built \val{val.bug.n} deliberately broken versions of the halo exchange and checked that the tests fail on them
(Table~\ref{tab:valbugs}). In the first version a received face is dropped, and in the second a rank sends, on one
face, the ghost layer it holds instead of its halo. Both change the parallel answer by a relative
\val{val.bug.min} to \val{val.bug.max}, and V3, V4 and V5 fail on both. V1 and V2 pass on both because a one-process run
has no halo between ranks, so a validation made only of serial checks would have accepted both errors.

\begin{table}[H]\centering\small
%% generated by scripts/report_macros.py; do not edit
\begin{tabular}{lccccc}
\toprule
broken halo exchange & V1 & V2 & V3 & V4 & V5 \\
\midrule
one received face is dropped & passes & passes & \textbf{fails} & \textbf{fails} & \textbf{fails} \\
the wrong plane is sent on one face & passes & passes & \textbf{fails} & \textbf{fails} & \textbf{fails} \\
\bottomrule
\end{tabular}

\caption{The two broken versions of the halo exchange against the five tests. A test that fails has detected the
error.}\label{tab:valbugs}
\end{table}

\emph{How accurate pipelined CG can be.} Outside the validation tests we asked both variants for a relative residual of
\val{val.floor.tol} on the 2D grid with $n=\val{val.floor.n}$ on \val{val.floor.ranks} ranks, allowing at most
\val{val.floor.maxit} iterations. Classic CG \val{val.floor.classic.verdict}. Pipelined CG
\val{val.floor.pipelined.verdict}. The smallest true residual that a variant can reach in floating point is called its
\term{attainable accuracy}, and it is worse for pipelined CG, a true residual of \val{val.floor.pipelined.true} against
\val{val.floor.classic.true} for classic CG in this run. The two variants agree at the tolerances used in the validation
tests (V5) and differ only when a residual far below those tolerances is requested.

% ======================================================================================================
\section{Hardware and measurement}\label{sec:env}\label{sec:timing}
\req{E10}
All timed runs were made at NUS HPC, on \val{cl.used}. A CPU node has \val{node.sockets} sockets, each an
\val{node.cpu} processor with \val{node.cps} cores, so \val{cpn} cores per node, and \val{env.mem}\,GB of memory in
total, divided between the two sockets. Intel's specification of this processor gives each socket
\val{spec.cpu.channels} memory channels of type \val{spec.cpu.memtype} at up to \val{spec.cpu.mts} million transfers per second, a peak of
\val{spec.cpu.peak}\,GB/s per socket and \val{spec.node.peak}\,GB/s per node, and \val{spec.cpu.cache}\,MB of cache
shared by the cores of a socket. We did not record how the memory modules are fitted, which decides whether that peak
rate applies, and Section~\ref{sec:mpi} compares it with the bandwidth we measured. A core reads the memory of its own
socket directly and that of the other socket through the link between the two. One core draws only a small part of a
socket's bandwidth, so the first few cores of a socket add bandwidth almost in proportion: two cores run the 3D
strong-scaling grid \val{pst.3.classic.2.S} times as fast as one, and four cores \val{pst.3.classic.4.S} times
(Table~\ref{tab:strong}). Once the cores of a socket together saturate its memory channels, the bandwidth available to a
run grows with the number of sockets it uses, not with its number of cores. The launcher
places the ranks of a node on consecutive cores and binds each rank to its cores, so the ranks fill the first socket
before the second: the binding report of a full-node run shows the first \val{node.cps} ranks on the first socket. A run
with at most \val{node.cps} ranks on a node therefore uses one socket and only that socket's memory bandwidth. A GPU node,
named \val{env.gpu.node} in the cluster, has the same processors and \val{env.gpus} NVIDIA GPUs. Our jobs did not record
the GPU model, which we take from the node name to be the NVIDIA \val{spec.gpu.name}. NVIDIA's specification of it gives
\val{spec.gpu.mem}\,GB of memory and a memory bandwidth of \val{spec.gpu.bw}\,GB/s, against the \val{spec.node.peak}\,GB/s
peak of the two sockets of a CPU node.

The CPU runs were compiled with GCC \val{env.cxx} and run with Open MPI \val{env.mpi}, and the GPU runs were compiled
with the NVIDIA HPC SDK compiler nvc++ \val{env.gpu.cxx} and run with the Open MPI \val{env.gpu.mpi} of its HPC-X
package, each on a compute node of the type that ran it. A laptop was used only to develop the program and to run the
validation tests during development, and none of its timings appear in this report.

Each timed run performs exactly \val{iters} CG iterations, whatever the residual, so that runs of different sizes and
layouts do the same work per unknown and the total cluster time stays bounded. The time of a solve is that of the
slowest rank, measured with \code{MPI\_Wtime} from a barrier at its start.

Each run solves the problem
\val{repeats} times in one launch of the program, and we report the median. Over the \val{rep.cells} configurations
solved \val{repeats} times, the largest and the smallest of their times differ by a median of \val{rep.median} of
the median time, and by more than \val{rep.thr} in \val{rep.nthr} of them. The largest differences are in one-core
runs, up to \val{rep.serial.max1s} for runs longer than a second, and the median keeps a single outlying repeat out of
the reported time. The runs to a tolerance (Section~\ref{sec:serial}) and the labelling runs of
Section~\ref{sec:graph}, which measure iteration counts and accuracy, are solved once, as are the one-core graph runs
on com-Orkut, which take about \val{gr.longrep} minutes per solve. The \val{rep.long.n} repeated runs of more than
\val{rep.long.min} minutes differ by at most \val{rep.long.max} between repeats.

The halo, reduction and computation times are averaged over ranks,
so the time a rank spends waiting inside a reduction for a slower rank counts as reduction time, and this
\term{load imbalance} is part of what we report as the cost of communication. Parallel runs used whole nodes reserved
for the job, so no other job shared a node's memory bandwidth or network links during a
measurement.\val{spread.except}\val{env.serialnote} Table~\ref{tab:plan} lists the experiments.

\begin{table}[H]\centering\small
%% generated by scripts/report_macros.py; do not edit
\begin{tabular}{>{\raggedright\arraybackslash}p{0.15\linewidth}>{\raggedright\arraybackslash}p{0.19\linewidth}>{\raggedright\arraybackslash}p{0.19\linewidth}lrl}
\toprule
experiment & what varies & problem & cores & configs. & solves \\
\midrule
serial, fixed iterations & grid size & 2D $n=64$--$8192$, 3D $n=16$--$406$ & 1 & 16 & 3 \\
serial, to a tolerance & grid size & 2D $n=64$--$1024$, 3D $n=16$--$101$ & 1 & 10 & 1 \\
strong scaling & cores, CG variant & 2D $n=8192$, 3D $n=406$ & 1--288 & 40 & 3 \\
weak scaling & cores and grid together & about $2.1\times10^{6}$ unknowns per rank & 1--288 & 40 & 3 \\
methods compared & decomposition, halo exchange, CG variant & strong-scaling grids & 72 and 288 & 40 & 3 \\
ranks $\times$ threads & ranks and threads per node & strong-scaling grids & 72 and 288 & 20 & 3 \\
GPU & number of GPUs & strong-scaling grids & 1 and 2 GPUs & 4 & 3 \\
graph, serial & graph, vertex order & all six graphs & 1 & 12 & 3 \\
graph, strong scaling & cores, CG variant & com-LiveJournal, com-Orkut & 1--288 & 36 & 3, 1 for 2 \\
graph, methods compared & partition, ghost exchange, CG variant, ranks $\times$ threads & com-LiveJournal, com-Orkut & 72 and 288 & 61 & 3 \\
graph, labelling & shift $a$, seeds $s$, edges kept & all six graphs & 8 and 72 & 70 & 1 \\
\bottomrule
\end{tabular}

\caption{The experiments on \val{cl.used}: what each varies, the problem, the cores, the number of configurations and
how many times each configuration is solved.}\label{tab:plan}
\end{table}

% ======================================================================================================
\section{Serial runtime scaling with problem size}\label{sec:serial}
\req{E2}\req{E14}
The serial runs, on one core of \val{cl.serial}, vary the matrix dimension and the number of nonzeros together: 2D grids
with \val{serial.list2} points per direction and 3D grids with \val{serial.list3}, from \val{serial.Nmin} to
\val{serial.Nmax} unknowns. We chose the 3D sizes so that each has about as many unknowns as a 2D size, while its matrix
has seven nonzeros per row against five, so comparing the two members of a pair shows whether the time follows the
unknowns or the nonzeros. The cost
model of Section~\ref{sec:model} predicts a time per iteration proportional to $N$, which is a slope of one in
Figure~\ref{fig:fig_serial}.

\figslot{fig_serial}{0.98\linewidth}{Serial time for \val{iters} iterations against the number of unknowns $N$ (left)
and the number of nonzeros (right), 2D and 3D grids, one core of \val{cl.serial}.}

\begin{table}[H]\centering\small
%% generated by scripts/report_macros.py; do not edit
\begin{tabular}{lrrrrrr}
\toprule
grid & $n$ & unknowns $N$ & nonzeros & time (s) & ns per unknown & ns per nonzero \\
\midrule
2D & 64 & 4{,}096 & 20{,}224 & 0.00279 & 6.82 & 1.38 \\
2D & 128 & 16{,}384 & 81{,}408 & 0.00426 & 2.60 & 0.52 \\
2D & 256 & 65{,}536 & 326{,}656 & 0.0209 & 3.19 & 0.64 \\
2D & 512 & 262{,}144 & 1{,}308{,}672 & 0.0978 & 3.73 & 0.75 \\
2D & 1{,}024 & 1{,}048{,}576 & 5{,}238{,}784 & 0.374 & 3.57 & 0.71 \\
2D & 2{,}048 & 4{,}194{,}304 & 20{,}963{,}328 & 1.72 & 4.09 & 0.82 \\
2D & 4{,}096 & 16{,}777{,}216 & 83{,}869{,}696 & 5.86 & 3.49 & 0.70 \\
2D & 8{,}192 & 67{,}108{,}864 & 335{,}511{,}552 & 24.1 & 3.59 & 0.72 \\
3D & 16 & 4{,}096 & 27{,}136 & 0.00392 & 9.58 & 1.45 \\
3D & 25 & 15{,}625 & 105{,}625 & 0.00401 & 2.56 & 0.38 \\
3D & 40 & 64{,}000 & 438{,}400 & 0.0164 & 2.56 & 0.37 \\
3D & 64 & 262{,}144 & 1{,}810{,}432 & 0.109 & 4.15 & 0.60 \\
3D & 101 & 1{,}030{,}301 & 7{,}150{,}901 & 0.404 & 3.92 & 0.57 \\
3D & 161 & 4{,}173{,}281 & 29{,}057{,}441 & 1.74 & 4.18 & 0.60 \\
3D & 256 & 16{,}777{,}216 & 117{,}047{,}296 & 6.22 & 3.71 & 0.53 \\
3D & 406 & 66{,}923{,}416 & 467{,}474{,}896 & 28.4 & 4.24 & 0.61 \\
\bottomrule
\end{tabular}

\caption{Serial runs on \val{cl.serial}: median time for \val{iters} iterations, and the cost of one iteration per
unknown and per nonzero.}\label{tab:serial}
\end{table}

\revised{serial}%
The cost per unknown is nearly constant once the vectors no longer fit in cache, as expected of a memory-bound
computation. The slopes fitted over all \val{ser.nsizes2} sizes of each dimension are \val{ser.slope2} in 2D and
\val{ser.slope3} in 3D, against the predicted value of one. From $N=\val{ser.2.n512.N}$ unknowns upwards, where the four
vectors that an iteration touches take \val{ser.ws.2.n512}\,MB or more, a 2D iteration costs between
\val{ser.2.n4096.ns} and \val{ser.2.n2048.ns}\,ns per unknown (Table~\ref{tab:serial}). The cost does not fall at
$N=\val{ser.2.n1024.N}$, whose vectors take \val{ser.ws.2.n1024}\,MB and would fit in the \val{spec.cpu.cache}\,MB cache
of a socket (Section~\ref{sec:env}), so for a single core that shared cache does not keep the iteration away from main
memory. The two smallest grids depart from this in opposite directions. The smallest costs \val{ser.2.n64.ns}\,ns per
unknown in 2D, because with only \val{ser.2.n64.N} unknowns the fixed cost of starting each loop is no longer small next
to the work, while the next one costs only \val{ser.2.n128.ns}\,ns, because its vectors, \val{ser.ws.2.n128}\,MB, fit
in the cache close to the core and the iteration does not wait for main memory. The serial repeats spread more than most
runs (Section~\ref{sec:timing}), so differences of a few per cent between neighbouring sizes are not significant.

The time follows the unknowns more than the nonzeros. On the largest grids a 3D iteration costs
\val{ser.ns3}\,ns per unknown against \val{ser.ns2}\,ns in 2D, a ratio of \val{ser.ratio32}, well below the ratio of
nonzeros per row of seven to five, and per nonzero the 3D iteration is the cheaper one, at \val{ser.nsnnz3} against
\val{ser.nsnnz2}\,ns. Most of the neighbour values that the stencil reads are still in cache from the previous plane or
line, so the time follows the vectors streamed from memory, which are the same number per unknown in 2D and 3D, rather
than the nonzeros.

The time to a fixed tolerance grows faster than $N$ because the iteration count grows with $n$
(Figure~\ref{fig:fig_serial_tol}). The fitted exponent of iterations against $n+1$ is \val{ser.itslope2} in 2D and
\val{ser.itslope3} in 3D, against the value of one that $\sqrt{\kappa}\propto n$ predicts. With a constant cost per
unknown this would make the time to the tolerance of \val{serialtol.tol} grow like $N^{1+1/d}$
(Section~\ref{sec:algorithm}), with exponents of \val{ser.tslope.pred2} in 2D and \val{ser.tslope.pred3} in 3D. The
fitted exponents of the time are larger, \val{ser.tslope2} and \val{ser.tslope3}, because the largest of these runs
has only \val{ser.2.n1024.N} unknowns (2D), in the range where the cost per unknown still rises as the vectors outgrow
the cache: fitted against $N$, the
cost per unknown and iteration of these runs grows with an exponent of \val{ser.tcslope2} in 2D and \val{ser.tcslope3}
in 3D, about the difference.

\figslot{fig_serial_tol}{0.98\linewidth}{Iterations to a relative residual of \val{serialtol.tol} against the grid
points per direction $n$ (left), and the serial time to reach it against $N$ (right), one core of \val{cl.serialtol}.}

% ======================================================================================================
\section{MPI strong and weak scaling, and the cost of communication}\label{sec:mpi}
\req{E3}\req{E4}\req{E5}\req{E14}
Strong scaling keeps the problem fixed and adds cores, and weak scaling grows the problem with the cores so that each
core keeps the same work. With $T(P)$ the time on $P$ cores, the speedup is $S=T(1)/T(P)$, the strong-scaling efficiency
is $E=S/P$ and the weak-scaling efficiency is $E=T(1)/T(P)$, so that $E=1$ is perfect. The serial time $T(1)$ is the
one-rank run of the same program and CG variant, compiled with MPI. The runs in this section were made on
\val{cl.strong} with one thread per rank. \val{scale.where}\val{scale.note}\val{spread.sent}

\subsection{Strong scaling}\label{sec:strong}
The strong-scaling grids have \val{strong.N2} unknowns in 2D ($n=\val{strong.n2}$) and \val{strong.N3} in 3D
($n=\val{strong.n3}$), so 2D and 3D have about the same number of unknowns, while the 3D matrix has seven nonzeros per
row against five and a 3D rank has more halo. They fit in one node's memory, at most \val{strong.mem3}\,GB for pipelined CG, and \val{iters} iterations of classic CG
take \val{st.2.classic.T1}\,s in 2D and \val{st.3.classic.T1}\,s in 3D on one core. We ran \val{st.counts} ranks. With
classic CG the efficiency at \val{meas.top} cores\val{meas.where} is \val{st.2.classic.Etop} in 2D and
\val{st.3.classic.Etop} in 3D, from speedups of \val{st.2.classic.Stop} and \val{st.3.classic.Stop} against the ideal
\val{meas.top}, and on one full node it is \val{st.2.classic.Enode} and \val{st.3.classic.Enode}. With pipelined CG and
the non-blocking halo exchange the efficiency at \val{meas.top} cores is \val{st.2.pipelined.Etop} in 2D and
\val{st.3.pipelined.Etop} in 3D.\val{meas.both}

\figslot{fig_strong}{0.98\linewidth}{Strong scaling on \val{cl.strong}: time for \val{iters} iterations on the fixed 2D
(left) and 3D (right) grids against cores, for classic CG with blocking halo exchange and pipelined CG with non-blocking
halo exchange. The dotted line is perfect scaling from the classic one-core time.\val{spread.cap}}

\begin{table}[H]\centering\small
%% generated by scripts/report_macros.py; do not edit
\begin{tabular}{rrrrrrr}
\toprule
cores & classic T (s) & $S$ & $E$ & pipelined T (s) & $S$ & $E$ \\
\midrule
1 & 24.1 & 1.0 & 1.00 & 26.5 & 1.0 & 1.00 \\
2 & 12.1 & 2.0 & 0.99 & 13.4 & 2.0 & 0.99 \\
4 & 6.62 & 3.6 & 0.91 & 7.40 & 3.6 & 0.90 \\
9 & 4.20 & 5.7 & 0.64 & 4.95 & 5.4 & 0.60 \\
18 & 3.14 & 7.7 & 0.43 & 4.12 & 6.4 & 0.36 \\
36 & 2.75 & 8.7 & 0.24 & 3.79 & 7.0 & 0.19 \\
72 & 1.36 & 17.8 & 0.25 & 1.89 & 14.1 & 0.20 \\
144 & 0.664 & 36.2 & 0.25 & 0.936 & 28.3 & 0.20 \\
\multicolumn{7}{c}{\pendingrow{vanda strong: 32 of 40 runs recorded}} \\
\bottomrule
\end{tabular}
\\[4pt]
%% generated by scripts/report_macros.py; do not edit
\begin{tabular}{rrrrrrr}
\toprule
cores & classic T (s) & $S$ & $E$ & pipelined T (s) & $S$ & $E$ \\
\midrule
1 & 28.4 & 1.0 & 1.00 & 30.5 & 1.0 & 1.00 \\
2 & 14.2 & 2.0 & 1.00 & 15.5 & 2.0 & 0.98 \\
4 & 7.35 & 3.9 & 0.96 & 8.05 & 3.8 & 0.95 \\
9 & 4.41 & 6.4 & 0.72 & 5.09 & 6.0 & 0.67 \\
18 & 3.43 & 8.3 & 0.46 & 4.48 & 6.8 & 0.38 \\
36 & 3.01 & 9.4 & 0.26 & 4.26 & 7.2 & 0.20 \\
72 & 1.50 & 19.0 & 0.26 & 2.16 & 14.1 & 0.20 \\
144 & 0.762 & 37.2 & 0.26 & 1.05 & 29.0 & 0.20 \\
288 & 0.377 & 75.3 & 0.26 & 0.542 & 56.3 & 0.20 \\
\bottomrule
\end{tabular}

\caption{Strong scaling on \val{cl.strong}, 2D (top) and 3D (bottom): median time $T$ for \val{iters} iterations,
speedup $S$ and efficiency $E$.\val{spread.tabnote}}\label{tab:strong}
\end{table}

\revised{mpi-node}%
Up to \val{node.cpn2} cores almost all of the efficiency is lost inside the first node. With classic CG on the 3D grid
the \val{node.cpn} cores of one node run \val{pst.3.classic.node.S} times as fast as one core (\val{pst.2.classic.node.S}
times in 2D), where perfect scaling would give \val{node.cpn}, so the efficiency on one node is
\val{pst.3.classic.node.E}. Going from one node to two divides the time by \val{pst.3.classic.gain.node.2node}, close to
the ideal factor of two, so the efficiency on \val{node.cpn2} cores stays at \val{pst.3.classic.2node.E}. This is the
memory regime of the cost model (Section~\ref{sec:model}). Every iteration of classic CG streams about \val{cg.bytes.classic} bytes per unknown through memory,
whatever the number of cores, so once the cores of a socket saturate its memory bandwidth, more cores on that socket
cannot make an iteration faster, while a second socket or a second node brings its own memory. We express this as an
effective bandwidth, the unknowns times the bytes moved per unknown and iteration (\val{cg.bytes.classic} for classic CG
and \val{cg.bytes.pipelined} for pipelined CG, from the memory model of Section~\ref{sec:model}) times the iterations,
divided by the measured time. Its median is \val{cal.core}\,GB/s over the large one-core runs and \val{cal.node}\,GB/s
over the full-node runs, a ratio of \val{cal.ratio}, and these are \val{cal.core.ofpeak} of the specified peak of one
socket and \val{cal.node.ofpeak} of the peak of a node (Section~\ref{sec:env}). The ratio comes from the same runs as the
speedup, so it restates the speedup in units of bandwidth and does not confirm it independently.

The shape of the curve inside the node follows from the two sockets described in Section~\ref{sec:env}. From
\val{node.q} to \val{node.half} cores the time falls only by a factor of \val{pst.3.classic.gain.q.half}, against the
ideal two, and from \val{node.half} to \val{node.cpn} cores it halves (\val{pst.3.classic.gain.half.node}). The launcher
fills the first socket before the second, so a run with \val{node.half} ranks uses one socket and only that socket's
memory channels, which were already close to saturation at \val{node.q} cores, and the second socket doubles the
bandwidth. The points between \val{node.q} and \val{node.half} cores therefore measure one socket, not half a node. We
did not rerun them with the ranks spread over both sockets, which would show how much higher their efficiency would be.

\subsection{Weak scaling}
Weak scaling keeps about \val{weak.per2} unknowns per rank in 2D (base $n=\val{weak.n2}$) and \val{weak.per3} in 3D
(base $n=\val{weak.n3}$), growing $n$ like $P^{1/d}$, so at \val{meas.top} ranks the grids have $n=\val{wk.top.n2}$ in 2D
and $n=\val{wk.top.n3}$ in 3D. The number of iterations is fixed, so this measures the cost of one iteration, while to a
fixed tolerance the iteration count would also grow with $n$ (Section~\ref{sec:serial}). With classic CG the efficiency
at \val{meas.top} cores\val{meas.where} is \val{wk.2.classic.Etop} in 2D and \val{wk.3.classic.Etop} in 3D, and on one
full node it is \val{wk.2.classic.Enode} and \val{wk.3.classic.Enode}.

Weak scaling shows the same memory limit from the other side. With the work per rank fixed, the 3D time rises from
\val{pwk.3.classic.one.T}\,s on one rank to \val{pwk.3.classic.half.T}\,s with \val{node.half} ranks sharing one socket,
stays at \val{pwk.3.classic.node.T}\,s on the full node, where each socket carries the same load, and grows by a factor
of only \val{pwk.3.classic.ratio.2node.node} on two nodes. If memory bandwidth were the only limit, a node's
\val{cal.node}\,GB/s shared by \val{node.cpn} ranks would give each rank \val{cal.share} of what one rank gets alone,
slightly more than the measured one-node efficiency of \val{pwk.3.classic.node.E} in 3D and \val{pwk.2.classic.node.E}
in 2D. Up to \val{node.cpn2} cores, therefore, communication between nodes costs little next to the memory limit inside
a node.

\figslot{fig_weak}{0.98\linewidth}{Weak scaling on \val{cl.strong}: time for \val{iters} iterations with the work per
core held fixed, 2D (left) and 3D (right). A flat line is perfect scaling.\val{spread.cap}}

\begin{table}[H]\centering\small
%% generated by scripts/report_macros.py; do not edit
\begin{tabular}{rrrrr}
\toprule
cores & classic T (s) & $E$ & pipelined T (s) & $E$ \\
\midrule
1 & 0.558 & 1.00 & 0.725 & 1.00 \\
2 & 0.851 & 0.66 & 0.799 & 0.91 \\
4 & 0.909 & 0.61 & 0.910 & 0.80 \\
9 & 1.19 & 0.47 & 1.38 & 0.53 \\
18 & 1.77 & 0.32 & 2.32 & 0.31 \\
36 & 3.11 & 0.18 & 4.27 & 0.17 \\
72 & 3.13 & 0.18 & 4.28 & 0.17 \\
144 & 3.22 & 0.17 & 4.33 & 0.17 \\
288 & 3.41 & 0.16 & 4.53 & 0.16 \\
\bottomrule
\end{tabular}
\\[4pt]
%% generated by scripts/report_macros.py; do not edit
\begin{tabular}{rrrrr}
\toprule
cores & classic T (s) & $E$ & pipelined T (s) & $E$ \\
\midrule
1 & 0.642 & 1.00 & 0.770 & 1.00 \\
2 & 0.948 & 0.68 & 0.987 & 0.78 \\
4 & 0.987 & 0.65 & 1.09 & 0.71 \\
9 & 1.25 & 0.51 & 1.44 & 0.53 \\
18 & 1.96 & 0.33 & 2.57 & 0.30 \\
36 & 3.40 & 0.19 & 4.78 & 0.16 \\
72 & 3.41 & 0.19 & 4.70 & 0.16 \\
144 & 3.77 & 0.17 & 4.67 & 0.16 \\
\multicolumn{5}{c}{\pendingrow{vanda weak: 32 of 40 runs recorded}} \\
\bottomrule
\end{tabular}

\caption{Weak scaling on \val{cl.strong}, 2D (top) and 3D (bottom): median time for \val{iters} iterations and
efficiency.\val{spread.tabnote}}\label{tab:weak}
\end{table}

\figslot{fig_efficiency}{0.98\linewidth}{Parallel efficiency against cores for both scaling experiments and both CG
variants on \val{cl.strong}. The dotted line is perfect efficiency. Points up to \val{cpn} cores are inside one
node\val{scale.beyond}.\val{spread.cap}}

\subsection{The cost of communication}
Figure~\ref{fig:fig_comm} splits the classic strong-scaling runs into computation, halo exchange and reductions per
iteration, and Table~\ref{tab:comm} gives the split at one node and at \val{meas.top} cores, together with the number of
messages and halo values per exchange. At \val{meas.top} cores\val{meas.where} the halo exchange takes \val{cm.3.halo}
of the 3D run time and \val{cm.2.halo} of the 2D one, and the reductions take \val{cm.3.red} and \val{cm.2.red}.

Up to \val{node.cpn2} cores communication is a small part of the time. With classic CG on one node the halo exchange
takes \val{pcm.3.classic.node.halo} of the 3D time and the reductions \val{pcm.3.classic.node.red}, and on two nodes the
shares are \val{pcm.3.classic.2node.halo} and \val{pcm.3.classic.2node.red}. In 2D, where a rank's halo is a line rather
than a surface, they are \val{pcm.2.classic.node.halo} and \val{pcm.2.classic.node.red} on one node. The reductions grow
faster than the halo exchange when the run leaves the node, because every reduction then waits for a message between
nodes, while the halo of a block shrinks with its faces, as the cost model of Section~\ref{sec:model} predicts. Pipelined
CG cuts the share of the reductions on two nodes to \val{pcm.3.pipelined.2node.red}, but up to \val{node.cpn2} cores it
is slower at every core count, taking \val{pst.3.pratio.node} times the classic time on one node and
\val{pst.3.pratio.2node} times on two (\val{pst.2.pratio.node} and \val{pst.2.pratio.2node} in 2D). The memory model of
Section~\ref{sec:model} gives it \val{cg.bytes.pipelined} bytes of memory traffic per unknown against
\val{cg.bytes.classic} for classic CG, a ratio of \val{cg.byteratio}. On full nodes, where memory bandwidth is saturated
(Section~\ref{sec:strong}), the measured time ratios are a little above that ratio. On one core, which draws only a small
part of a socket's bandwidth, the ratio is only \val{pst.3.pratio.one} in 3D and \val{pst.2.pratio.one} in 2D
(Table~\ref{tab:strong}), so there the time is not set by the bytes moved, most likely by how many memory requests one
core can keep in flight, which we did not measure. Where the bandwidth is saturated, the extra traffic costs more than
the reduction time it hides as long as the reductions take only a small share of the time.

\figslot{fig_comm}{0.98\linewidth}{Time per iteration in computation, halo exchange and reductions, mean over ranks,
for classic CG with blocking halo exchange on \val{cl.strong}. At one rank there is no halo exchange or reduction to
time.\val{spread.cap}}

\begin{table}[H]\centering\footnotesize
%% generated by scripts/report_macros.py; do not edit
\begin{tabular}{lrrrrrrr}
\toprule
grid & cores & computation & halo exchange & reductions & communication share & messages & halo values \\
\midrule
2D & 72 & 13.165 & 0.127 & 0.231 & 2.6\,\% & 4 & 4{,}098 \\
2D & 288 & 2.665 & 0.096 & 0.568 & 19.9\,\% & 4 & 2{,}050 \\
3D & 72 & 13.503 & 0.671 & 0.793 & 9.8\,\% & 6 & 59{,}772 \\
3D & 288 & 2.733 & 0.467 & 0.563 & 27.4\,\% & 6 & 23{,}120 \\
\bottomrule
\end{tabular}

\caption{Time per iteration (ms) by part on \val{cl.strong}, classic CG with blocking halo exchange, at one node and at
\val{meas.top} cores, with the messages and halo values per exchange of the rank with the most.\val{spread.tabnote}}
\label{tab:comm}
\end{table}

\subsection{The methods compared}
Table~\ref{tab:variants} gives the time per iteration for every combination of decomposition, halo exchange and CG
variant on one node and on \val{topcores} cores. On one node the fastest 3D combination is the \val{var.3.node.best},
and the fastest 2D one is the \val{var.2.node.best}. At \val{topcores} cores the fastest 3D combination is the
\val{var.3.top.best} and the slowest takes \val{var.3.top.spread} times as long, and in 2D the fastest is the
\val{var.2.top.best} and the slowest takes \val{var.2.top.spread} times as long.

\begin{table}[H]\centering\footnotesize
%% generated by scripts/report_macros.py; do not edit
\begin{tabular}{llllrr}
\toprule
grid & decomposition & CG & halo exchange & ms/it, 72 cores & ms/it, 288 cores \\
\midrule
2D & 1D & classic & blocking & 13.585 & 3.192 \\
2D & 1D & classic & non-blocking & 13.575 & 3.543 \\
2D & 1D & pipelined & blocking & 18.788 & 4.849 \\
2D & 1D & pipelined & non-blocking & 18.814 & 4.806 \\
2D & 2D & classic & blocking & 13.560 & 3.343 \\
2D & 2D & classic & non-blocking & 13.582 & 3.207 \\
2D & 2D & pipelined & blocking & 18.691 & 4.713 \\
2D & 2D & pipelined & non-blocking & 18.875 & 4.668 \\
3D & 1D & classic & blocking & 19.714 & 8.990 \\
3D & 1D & classic & non-blocking & 19.783 & 8.951 \\
3D & 1D & pipelined & blocking & 25.358 & 10.734 \\
3D & 1D & pipelined & non-blocking & 25.292 & 10.710 \\
3D & 2D & classic & blocking & 14.576 & 3.802 \\
3D & 2D & classic & non-blocking & 14.668 & 3.659 \\
3D & 2D & pipelined & blocking & 19.718 & 5.284 \\
3D & 2D & pipelined & non-blocking & 19.851 & 5.149 \\
3D & 3D & classic & blocking & 14.967 & 3.767 \\
3D & 3D & classic & non-blocking & 15.607 & 4.048 \\
3D & 3D & pipelined & blocking & 20.399 & 5.313 \\
3D & 3D & pipelined & non-blocking & 21.560 & 5.421 \\
\bottomrule
\end{tabular}

\caption{Time per iteration (ms) for every combination of decomposition, CG variant and halo exchange, at
\val{plan.variantsP} cores on \val{cl.variants}.}\label{tab:variants}
\end{table}

\interpslot{mpi-top}{From the runs at \val{topcores} cores: whether pipelined CG overtakes classic CG at \val{topcores} cores, where
the reductions take a larger share, whether the non-blocking halo exchange hid any message time, and how the slab
decomposition behaves at \val{topcores} ranks (Table~\ref{tab:geometry}).}

% ======================================================================================================
\section{MPI with OpenMP threads, and the GPU}\label{sec:mpix}
\req{E6}\req{E7}
Lecture~9 asks four questions of an MPI+OpenMP program, and we answer them in its order. The first is what is
parallelised across ranks and what with threads. Ranks own boxes of the grid, so the domain decomposition, the halo
exchange and the reductions belong to MPI, while threads split every loop over a rank's own box: the stencil, the vector
updates and the inner products (Section~\ref{sec:design}).

The second question is how ranks and threads synchronise and how much synchronisation is needed. Each threaded loop ends
at OpenMP's implicit barrier, after which the master thread alone makes the MPI calls: the halo exchange before the
stencil and one reduction for each inner product, whose local part the threads have already summed. This needs one
barrier per loop and no locks, and it is why the library runs at the \code{MPI\_THREAD\_FUNNELED} level. The third
question is how threads interact with MPI windows, and in our program they do not, because the halo exchange uses
two-sided messages and no windows (Section~\ref{sec:design}).

The fourth question is which split of a \val{lay.cpn}-core node into ranks and threads is fastest. We timed
\val{plan.layouts} ranks $\times$ threads per node, on \val{lay.nodes} of \val{cl.layout}, with the strong-scaling grids.
Fewer ranks mean fewer and larger messages and a shallower reduction, but threads pay for starting and joining every
parallel loop and for memory placed across sockets. On one node the fastest layout is \val{lay.3.1.best} in 3D and
\val{lay.2.1.best} in 2D\val{lay.top.sent}. On \val{lay.topnodes.phrase} the slowest 3D layout takes
\val{lay.3.top.spread} times as long as the fastest, and the slowest 2D layout \val{lay.2.top.spread} times.

\figslot{fig_layout}{0.98\linewidth}{Time for \val{iters} iterations for each split of a node's \val{cpn} cores into
ranks and threads, on \val{lay.nodes} of \val{cl.layout}. The right end of each line is pure MPI.}

\begin{table}[H]\centering\small
%% generated by scripts/report_macros.py; do not edit
\begin{tabular}{lrr}
\toprule
ranks $\times$ threads per node & 2D, 1 node & 3D, 1 node \\
\midrule
1$\times$72 & 1.47 & 1.81 \\
2$\times$36 & 1.35 & 1.59 \\
4$\times$18 & 1.34 & 1.48 \\
8$\times$9 & 1.35 & 1.47 \\
72$\times$1 & 1.36 & 1.50 \\
\multicolumn{3}{c}{\pendingrow{vanda layout: 10 of 20 runs recorded}} \\
\bottomrule
\end{tabular}

\caption{Median time (s) for \val{iters} iterations by ranks $\times$ threads per node, on \val{cl.layout}.}
\label{tab:layout}
\end{table}

The GPU runs use OpenMP GPU offload on a \val{cl.gpu} GPU node, one rank per GPU, with the strong-scaling grids. One GPU
takes \val{gpu.3.t}\,s for the 3D grid and \val{gpu.2.t}\,s for the 2D grid, which is \val{gpu.3.vsnode} and
\val{gpu.2.vsnode} times as fast as the best layout on one CPU node (Table~\ref{tab:gpu}). In the two-GPU runs the halos
also travel from each GPU to its host and back.

\begin{table}[H]\centering\small
%% generated by scripts/report_macros.py; do not edit
\begin{tabular}{lrrrr}
\toprule
grid & GPUs & GPU time (s) & fastest 1-node CPU time (s) & CPU node / GPU \\
\midrule
2D & 1 & 1.20 & 1.34 & 1.12 \\
2D & 2 & 0.611 & 1.34 & 2.20 \\
3D & 1 & 1.26 & 1.46 & 1.16 \\
3D & 2 & 0.712 & 1.46 & 2.05 \\
\bottomrule
\end{tabular}

\caption{GPU runs against the best one-node CPU layout, on the same grid for \val{iters} iterations.}\label{tab:gpu}
\end{table}

\interpslot{mpix}{Explain which layout won and why (messages, reduction depth, memory placed across sockets, thread
start-up per loop), and whether the answer changed between one node and several, and explain the GPU result in terms of
memory bandwidth and what the two-GPU run says about the host copies of halos.}

% ======================================================================================================
\section{Real data: community labelling on SNAP graphs}\label{sec:graph}
\req{E11}\req{E12}\req{E14}
The project asks for an application to real data with conclusions about the data, and for a graph it asks how sparsity
affects the results. A graph also gives the solver an irregular communication pattern, unlike the grid, where each rank
exchanges a thin halo with a few neighbouring blocks. We therefore apply the same CG code to \val{gr.ndata} graphs from the Stanford Large Network Dataset Collection, which we call SNAP below:
\val{gr.datalist}. Each graph comes with known groups of vertices, its \term{communities}. We ask a semi-supervised question of the data: if a few vertices of each community are labelled,
how many of the others does diffusion over the graph label correctly, and how does that depend on how sparse the graph
is? The answer can be scored because the communities are known.

\subsection{The data and the communities}\label{sec:snapdata}
In email-Eu-core an edge records e-mail between two members of a European research institution, and SNAP gives the
department of each member in a separate file, so a community is a department. In the five \emph{com-} graphs, taken
from SNAP's networks with ground-truth communities, an edge joins two co-purchased products (Amazon), two co-authors
(DBLP) or two friends (Youtube, LiveJournal, Orkut), and a community is a product category, a publication venue or a
user group. For these five SNAP also provides a file (\code{.top5000.cmty}) of the communities that rank best, on
average, under several structural scores, among them conductance, which measures how many of a community's edges leave
it~\cite{yl}, and we use that file. Table~\ref{tab:graphdata} lists the graphs.

We downloaded the files from SNAP and converted them into the solver's input format, in which every graph is undirected
and simple: email-Eu-core is directed, so an edge in either direction becomes one undirected edge, and self-loops and
repeated edges are dropped. The vertices are numbered $0,\dots,V-1$ in ascending SNAP id, and the communities are mapped
to the same numbers. The converted vertex and edge counts can be compared with SNAP's published counts in the table,
where SNAP counts the directed edges of email-Eu-core.

\begin{table}[H]\centering\small
%% generated by scripts/report_macros.py; do not edit
\begin{tabular}{lrrrrrrr}
\toprule
dataset & $V$ & $E$ & SNAP $V$ / $E$ & $2E/V$ & $d_{\max}$ & evaluated & majority \\
\midrule
email-Eu-core & 1{,}005 & 16{,}064 & 1{,}005 / 25{,}571 & 32.0 & 345 & 538 & 0.193 \\
com-Amazon & 334{,}863 & 925{,}872 & 334{,}863 / 925{,}872 & 5.5 & 549 & 1{,}686 & 0.192 \\
com-DBLP & 317{,}080 & 1{,}049{,}866 & 317{,}080 / 1{,}049{,}866 & 6.6 & 343 & 31{,}391 & 0.227 \\
com-Youtube & 1{,}134{,}890 & 2{,}987{,}624 & 1{,}134{,}890 / 2{,}987{,}624 & 5.3 & 28{,}754 & 6{,}626 & 0.322 \\
com-LiveJournal & 3{,}997{,}962 & 34{,}681{,}189 & 3{,}997{,}962 / 34{,}681{,}189 & 17.3 & 14{,}815 & 8{,}178 & 0.173 \\
com-Orkut & 3{,}072{,}441 & 117{,}185{,}083 & 3{,}072{,}441 / 117{,}185{,}083 & 76.3 & 33{,}313 & 37{,}632 & 0.122 \\
\bottomrule
\end{tabular}

\caption{The graphs as converted: vertices $V$, undirected edges $E$, average degree $2E/V$ and largest degree
$d_{\max}$, SNAP's published counts, the number of evaluated vertices at $s=\val{gr.s}$ seeds per community and the
share of the largest of the \val{gr.K} communities among them.}\label{tab:graphdata}
\end{table}

We use the $K=\val{gr.K}$ largest communities of each graph that are nearly disjoint. Going through the communities
from the largest down, we take one if at most a fraction $\theta=\val{gr.theta}$ of its members belong to communities
already taken, and if it holds at most that fraction of the members of each of them. The largest communities leave
enough members both to seed and to evaluate, and near-disjointness matters because the top communities overlap heavily
and some are near-duplicates of each other, so that with two copies of the same group a correct label could not be told
from a wrong one. A vertex that belongs to exactly one of the $K$ communities we call \term{exclusive}. The seeds are
chosen among the exclusive vertices, and the \term{evaluated vertices} are the exclusive vertices that are not seeds,
so their number falls slightly as the number of seeds grows.

\subsection{The method}\label{sec:graphmethod}
Let $A$ be the adjacency matrix, $D$ the diagonal matrix of degrees and $L=D-A$ the graph Laplacian. For each community
$c$ we choose $s$ of its exclusive vertices as its \term{seed vertices} $S_c$ and, with a positive \term{shift} $a$,
solve
\begin{equation}
(L+aI)\,x_c=b_c,\qquad (b_c)_v=\begin{cases}1 & v\in S_c\\ 0 & \text{otherwise,}\end{cases}\label{eq:graph}
\end{equation}
then label every evaluated vertex $v$ with the community $c$ whose $(x_c)_v$ is largest. Row $v$ of \eqref{eq:graph}
reads $(x_c)_v=\big((b_c)_v+\sum_{u\sim v}(x_c)_u\big)/(d_v+a)$: away from the seeds each value is the sum of its
neighbours' values divided by $d_v+a$, so $x_c$ is a diffusion from the seeds damped by a factor of about $d_v/(d_v+a)$
per step, and a shift that is small against the degrees lets the labels travel far while a large one keeps them near
the seeds. Since $L$ is symmetric positive semidefinite and $a>0$, the matrix $L+aI$ is symmetric positive definite with
eigenvalues in $[a,\lambda_{\max}+a]$, where $\lambda_{\max}\le 2d_{\max}$ by Gershgorin's theorem, so its condition
number is at most $(2d_{\max}+a)/a$ and CG needs more iterations, roughly as $a^{-1/2}$, as $a$ shrinks. The seeds of
each community are its $s$ exclusive members with the smallest hash of their SNAP id, so they are the same for every
vertex order and rank count, the seeds for a smaller $s$ are the first of those for a larger one, and the seed lists are
recorded with the runs.

A vertex whose $x_c$ is zero for every $c$ is \term{unreached} and counted as wrong. This happens when no seed lies in
its connected component, and also because of a property of CG: after $k$ iterations $x_c$ lies in
$\mathrm{span}\{b_c,(L+aI)b_c,\dots,(L+aI)^{k-1}b_c\}$, which is zero beyond graph distance $k-1$ from the seeds. A
vertex is therefore labelled only if CG ran for more iterations than its distance from the nearest seed, so a large
$a$, which needs few iterations, also limits the reach. The accuracy is the share of evaluated vertices labelled
with their own community, and we compare it with two baselines, \emph{majority}, which always answers the largest
community among the evaluated vertices, and \emph{random}, whose expected accuracy is $1/K$. Unless stated otherwise we
use $a=\val{gr.a}$ and $s=\val{gr.s}$, and every system is solved to a relative residual of \val{gr.tol}.

\subsection{Parallel design for graphs}
The graph solver reuses the CG code of the Poisson problem unchanged and replaces only the operator. Both CG variants
are written once for any operator, and the stencil and the graph each supply a product $y=Ax$ that also returns the
local part of $x\cdot y$. The matrix $L+aI$ is stored in compressed sparse rows and partitioned in one dimension, each
rank owning a contiguous range of vertex numbers. The ranges are chosen in one of two ways, equal rows (\emph{rows})
giving each rank the same number of vertices, and equal nonzeros (\emph{nnz}) giving each rank the same number of
nonzeros of $L+aI$, which is the degree plus one per row. The degrees are heavy-tailed: a few vertices, which we call
\term{hubs}, have degrees far above the average ($d_{\max}$ against $2E/V$ in Table~\ref{tab:graphdata}), so equal rows
can give one rank far more work. We also number the vertices in a second
order, the breadth-first (Cuthill--McKee) order, which gives neighbours nearby numbers so that a contiguous range should
cut fewer edges.

The neighbours of a rank's vertices that are owned by other ranks are its \term{ghosts}, the graph's counterpart of the
halo, and their values must arrive before every product. The exchange lists are built once: each rank sends the ids it
needs to their owners, which then know what to send. Every product then exchanges the values in one of three ways:
\begin{itemize}
\item \code{MPI\_Alltoallv} over all ranks, with zero counts for ranks that share no edge,
\item \code{MPI\_Neighbor\_alltoallv} on a distributed graph communicator whose neighbours are only the ranks that do
share edges, which we call the \term{neighbourhood collective},
\item its non-blocking form \code{MPI\_Ineighbor\_alltoallv}, overlapped with the rows that have no ghost neighbour.
\end{itemize}
The received values land directly in the ghost part of each vector, behind the owned part, so the vector kernels of the
Poisson code run unchanged on the owned part. Within a rank the rows are split among threads by equal numbers of
nonzeros. Every run records the \term{edge cut} (the share of edges whose two ends are owned by different ranks), the
ghosts and neighbours per rank, the \term{nonzero imbalance} (the largest number of nonzeros per rank over the mean) and
the time spent exchanging ghosts, as the halo time was recorded for the Poisson problem.

Weak scaling stays on the synthetic problem because a real graph has a fixed size. Growing the work with the cores
would need a family of graphs of the same structure, which only a synthetic generator could provide, and the six graphs
differ in their degree distribution as well as in size, whereas the Poisson grid grows with the cores exactly.

\subsection{Validation of the graph solver}\label{sec:graphval}
The graph solver has two validation tests of its own, G1 and G2, and a third check, G3, that both fail on a
deliberately broken version of the ghost exchange, as for the Poisson solver.
\begin{itemize}
\item \textbf{G1} runs email-Eu-core on two to five ranks, with threads, both partitions, all three exchanges and both
vertex orders, in \val{gr.val.layouts} configurations. The product $(L+aI)v$ equals the one-rank product to
\val{gr.val.spmv}, against a tolerance of \val{gr.val.spmvtol}, and after \val{gr.val.iters} iterations of classic or
pipelined CG the iterates agree to \val{gr.val.x}, against \val{gr.val.xtol}. The comparison stops after
\val{gr.val.iters} iterations because beyond that the rounding differences between two correct runs grow quickly on
this graph. They grow in the same way between one and two threads on one rank, where no ghost is exchanged, so a longer
comparison would measure rounding rather than correctness. The growth that we measured during development
is listed with G1 in the validation program (Appendix~\ref{app:code}). Run to convergence, the answers of the parallel
runs, classic and pipelined, differ from the one-rank classic answer by at most \val{gr.val.conv}, below the smallest of
their condition-number bounds, \val{gr.val.bound}, computed as in V5.
\item \textbf{G2} builds a small graph by hand: two cliques of five vertices, a bridge, a tail, a self-loop, a repeated
edge and a vertex whose only edge is a self-loop. The conversion must report what was built, \val{gr.val.tinyV} vertices
and \val{gr.val.tinyE} edges, and on one to three ranks $x$ must match a dense solve of the same system to
\val{gr.val.dense}, against a tolerance of \val{gr.val.densetol}.
\item \textbf{G3} checks the tests themselves. In a deliberately broken version the ghosts from each rank's first
neighbour are received into a scratch area and never reach the operator, and it fails \val{gr.val.bugby}: its
products are off by \val{gr.val.bugspmv}, and its labelling accuracy drops from
\val{gr.val.acc} to \val{gr.val.bugacc}.
\end{itemize}
These results were obtained on \val{gr.val.where}, and the program used for the graph runs passed G1 and G2 on the
cluster before any graph run.

\subsection{Accuracy against the baselines}\label{sec:graphacc}
\begin{table}[H]\centering\small
%% generated by scripts/report_macros.py; do not edit
\begin{tabular}{lrrrrrr}
\toprule
dataset & accuracy & unreached & majority & random & iterations/solve & time/solve (s) \\
\midrule
email-Eu-core & 0.758 & 0.030 & 0.193 & 0.100 & 102 & 0.00284 \\
com-Amazon & 1.000 & 0.000 & 0.192 & 0.100 & 113 & 0.0309 \\
com-DBLP & 0.651 & 0.000 & 0.227 & 0.100 & 136 & 0.0377 \\
com-Youtube & 0.886 & 0.000 & 0.322 & 0.100 & 708 & 0.504 \\
com-LiveJournal & 0.986 & 0.000 & 0.173 & 0.100 & 624 & 5.05 \\
com-Orkut & 0.835 & 0.000 & 0.122 & 0.100 & 1{,}024 & 27.4 \\
\bottomrule
\end{tabular}

\caption{Labelling with $a=\val{gr.a}$ and $s=\val{gr.s}$ on one node of \val{cl.gr.sweep}, \val{gr.sweep.ranks} ranks
(\val{gr.sweep.ranks.email} for email-Eu-core): accuracy over the evaluated vertices, the share left unreached, the two
baselines, CG iterations per system and time per system.}
\label{tab:graphacc}
\end{table}

\revised{graph-accuracy}%
With \val{gr.s} seeds per community every graph is labelled far better than both baselines (Table~\ref{tab:graphacc}).
The accuracy lies between \val{gr.acc.lo} and \val{gr.acc.hi}, against majority shares between \val{gr.maj.com-orkut}
and \val{gr.maj.com-youtube} and an expected \val{gr.rand} for a random guess. com-LiveJournal reaches
\val{gr.acc.com-lj}, com-Youtube \val{gr.acc.com-youtube} and com-Orkut \val{gr.acc.com-orkut}, and email-Eu-core, at
\val{gr.acc.email-Eu-core}, is the second hardest after com-DBLP. Diffusion labels a vertex correctly when more of its
paths lead to the seeds of its own community than to those of others, so the accuracy depends on how many of a
community's edges stay inside it, which we did not measure. How the communities were chosen gives one hint: those of
the com- graphs are SNAP's selection of the communities with the best structural scores (Section~\ref{sec:snapdata}),
while the departments of email-Eu-core were not selected in this way. That hint does not explain com-DBLP, whose
communities were selected like com-Amazon's and are still the hardest to recover.

More seeds help on most graphs. From $s=1$ to $s=\val{gr.smax}$, email-Eu-core rises from \val{gw.email-Eu-core.s1.acc}
to \val{gw.email-Eu-core.s20.acc}, com-Orkut from \val{gw.com-orkut.s1.acc} to \val{gw.com-orkut.s20.acc} and com-DBLP
from \val{gw.com-dblp.s1.acc} to \val{gw.com-dblp.s20.acc}, while com-Youtube gains little (\val{gw.com-youtube.s1.acc}
to \val{gw.com-youtube.s20.acc}) and com-Amazon is already at \val{gw.com-amazon.s1.acc} with a single seed.
com-LiveJournal shows the largest single step, from \val{gw.com-lj.s1.acc} with one seed to \val{gw.com-lj.s2.acc} with
two. A plausible reading is that a single seed can sit at the edge of its community and a second one covers the rest,
and the distance from each evaluated vertex to its nearest seed, which we did not record, would test it.

\subsection{The shift, the seeds and sparsity}
\figslot{fig_graph_sweep}{0.98\linewidth}{Accuracy (top) and CG iterations per system (bottom) against the shift $a$,
the seeds per community $s$ and the average degree, on \val{cl.gr.sweep}. Each sweep varies one parameter from
$a=\val{gr.a}$, $s=\val{gr.s}$ and all edges, and the average degree is lowered by keeping each edge only with
probability \val{gr.keeplist}.}

Figure~\ref{fig:fig_graph_sweep} varies one parameter at a time: $a\in\{\val{gr.avals}\}$, $s\in\{\val{gr.svals}\}$ and
the sparsity. The sparsity is changed by keeping each edge with probability $p\in\{\val{gr.keepvals}\}$, decided from the
two SNAP ids so that every rank count and vertex order sees the same sparser graph, while the communities and seeds stay
the same, so that a sparser graph changes only the edges. Comparing the six graphs mixes the effect of sparsity with that
of the graph, and this controlled thinning separates the two.

\revised{graph-sparsity}%
Thinning harms the sparse graphs far more than the dense ones. Keeping each edge with probability $1/8$ lowers the
average degree of com-Amazon from \val{gw.com-amazon.deg} to \val{gr.keep.com-amazon.deg}, and then
\val{gr.keep.com-amazon.unr} of its evaluated vertices are unreached (Section~\ref{sec:graphmethod}) and its accuracy
drops to \val{gr.keep.com-amazon.acc}. com-DBLP and com-Youtube fall to average degrees of \val{gr.keep.com-dblp.deg} and
\val{gr.keep.com-youtube.deg}, to accuracies of \val{gr.keep.com-dblp.acc} and \val{gr.keep.com-youtube.acc}, with
\val{gr.keep.com-dblp.unr} and \val{gr.keep.com-youtube.unr} of their evaluated vertices unreached. These degrees are
below one, the average degree under which an Erd\H{o}s--R\'enyi random graph no longer has a giant component, one that
holds a fixed share of all vertices, although com-Youtube still reaches most of its evaluated vertices, so the
threshold is only a rough guide for these graphs. On these three sparse graphs most of the accuracy lost is unreached
vertices: com-Amazon loses \val{gr.keep.com-amazon.loss}, com-DBLP \val{gr.keep.com-dblp.loss} and com-Youtube
\val{gr.keep.com-youtube.loss}. The same thinning lowers the average degree of com-Orkut from \val{gw.com-orkut.deg} to
\val{gr.keep.com-orkut.deg}, leaves only \val{gr.keep.com-orkut.unr} of its evaluated vertices unreached and costs it
\val{gr.keep.com-orkut.loss} of accuracy, down to \val{gr.keep.com-orkut.acc}, so on this dense graph most of the smaller loss is wrong labels on reached
vertices. A sparser graph is also cheaper to solve: CG
needs \val{gr.keep.com-amazon.itratio} times fewer iterations per system on the thinned com-Amazon and
\val{gr.keep.com-orkut.itratio} times fewer on the thinned com-Orkut, because fewer edges lower the largest degree and
with it the largest eigenvalue and the condition number.

The iterations follow the condition-number estimate in the shift: their fitted exponent against $a$ is
$\val{gr.aslope}$, against the $-1/2$ that $\sqrt{\kappa}$ predicts. Over most of the range the accuracy hardly depends on
$a$, so a small shift costs iterations without buying accuracy: com-Youtube needs \val{gw.com-youtube.amin.it}
iterations per system at $a=\val{gw.com-youtube.amin}$ for an accuracy of \val{gw.com-youtube.amin.acc}, against
\val{gw.com-youtube.it} at $a=1$ for the accuracy in Table~\ref{tab:graphacc}. Two graphs depart from this. com-DBLP
labels better with a small shift, \val{gw.com-dblp.amin.acc} at $a=\val{gw.com-dblp.amin}$ against
\val{gw.com-dblp.a100.acc} at $a=\val{gr.amax}$, which suggests that its communities are spread out so that labels must
travel far, although we did not measure the distances within its communities. com-Amazon, at the largest shift, leaves
\val{gw.com-amazon.a100.unr} of its vertices unreached, because CG then stops after \val{gw.com-amazon.a100.it}
iterations per system, and after $k$ iterations only the vertices within distance $k-1$ of a seed can have a non-zero
value (Section~\ref{sec:graphmethod}).

\subsection{Serial cost per iteration across the six graphs}
\figslot{fig_graph_serial}{0.98\linewidth}{Time per CG iteration on one core of \val{cl.gr.serial} against the nonzeros
(left) and the vertices (right) for the six graphs, in the natural and the breadth-first vertex order, \val{gr.siters}
iterations on one system per point.}

\begin{table}[H]\centering\small
%% generated by scripts/report_macros.py; do not edit
\begin{tabular}{lrrrrr}
\toprule
graph & $V$ & nonzeros & ns/nonzero, natural & ns/nonzero, BFS & BFS speed-up \\
\midrule
email-Eu-core & 1{,}005 & 33{,}133 & 1.56 & 1.19 & 1.31 \\
com-Amazon & 334{,}863 & 2{,}186{,}607 & 3.20 & 3.09 & 1.03 \\
com-DBLP & 317{,}080 & 2{,}416{,}812 & 3.27 & 2.68 & 1.22 \\
com-Youtube & 1{,}134{,}890 & 7{,}110{,}138 & 2.79 & 1.59 & 1.75 \\
com-LiveJournal & 3{,}997{,}962 & 73{,}360{,}340 & 2.11 & 2.19 & 0.96 \\
com-Orkut & 3{,}072{,}441 & 237{,}442{,}607 & 2.42 & 2.24 & 1.08 \\
\bottomrule
\end{tabular}

\caption{Serial cost of one CG iteration per nonzero for each graph on one core of \val{cl.gr.serial}, in the natural and
the breadth-first (BFS) order, and the speed-up that the BFS order gives.}\label{tab:graphserial}
\end{table}

\revised{graph-serial}%
The six graphs vary the vertices and the edges together (Figure~\ref{fig:fig_graph_serial}). Unlike the grid, whose
time follows the unknowns (Section~\ref{sec:serial}), a graph iteration costs about the same time per nonzero: the time
per iteration grows with the nonzeros with a fitted exponent of \val{gr.ser.slope} in the natural order and
\val{gr.ser.slopebfs} in the breadth-first order, against the value of one for a constant cost per nonzero. On
com-Orkut, the largest graph, one iteration costs \val{gr.ser.ns.natural}\,ns per nonzero in the natural order and
\val{gr.ser.ns.bfs}\,ns in the breadth-first order, which is \val{gs.vsgrid} times the \val{ser.nsnnz2}\,ns of the 2D
stencil, although per nonzero the graph streams fewer bytes. The operator stores no matrix values, because every
off-diagonal entry of $L+aI$ is $-1$, only a column index of \val{gs.bytes.nnz} bytes per entry and \val{gs.bytes.row}
bytes per row for the row pointer and the degree. With the CG vectors an iteration on com-Orkut therefore streams about
\val{gs.com-orkut.bpnnz} bytes per nonzero, against \val{ser.bpnnz2} for the 2D stencil, and at the one-core bandwidth
of Section~\ref{sec:strong} that takes about \val{gs.com-orkut.streamns}\,ns per nonzero, a small part of the measured
time. The rest goes to reading each neighbour's value, one read per nonzero at an address that only the column index
decides, which the cache can serve only when neighbours have nearby numbers, whereas the stencil finds its neighbours at
fixed offsets, mostly in cache. The smallest graph, email-Eu-core, fits in cache and costs \val{gs.email-Eu-core.ns}\,ns
per nonzero (Table~\ref{tab:graphserial}).

The breadth-first numbering speeds up the serial iteration on most of the graphs (Table~\ref{tab:graphserial}),
because the values of a row's neighbours are then read from nearby memory. The clearest case is com-DBLP, which runs
\val{gs.com-dblp.bfsratio} times as fast, with natural-order repeats that agree within \val{gs.com-dblp.natspread}. The
largest ratio, \val{gr.ser.bfsgain}, is uncertain: the natural-order times of com-Youtube fell from one repeat to the
next and spread over \val{gs.com-youtube.natspread} of their median, and with the fastest repeat of each order the
speed-up is \val{gs.com-youtube.bfsratio.fast}. The numbering helps least on the two largest graphs, which are also the
two densest com- graphs (Table~\ref{tab:graphdata}), where com-Orkut runs \val{gs.com-orkut.bfsratio} times as fast and com-LiveJournal \val{gs.com-lj.bfsratio} times,
slightly slower. In a dense graph a breadth-first search reaches a large share of the vertices within a few steps, so
each level of the search is large and the order cannot keep all of a vertex's neighbours close to it, and the parallel
runs below show the same effect as a larger edge cut. We did not record the level sizes that would confirm this.
com-Amazon, which is sparse, also gains little (Table~\ref{tab:graphserial}), for a reason we did not establish.

\subsection{Strong scaling and communication}
\figslot{fig_graph_strong}{0.98\linewidth}{Strong scaling on \val{cl.gr.strong}: time per system against cores for
\val{gr.strongdata}, classic and pipelined CG. Each point times \val{gr.rhs} of the \val{gr.K} systems, and the dotted
line is perfect scaling from one core.\val{spread.cap}}

The two largest graphs, \val{gr.strongdata}, were run on \val{gr.strong.counts} cores, one rank per core, with equal
nonzeros, the neighbourhood collective and the natural order. At \val{gr.st.top} cores classic CG reaches a speedup of
\val{gr.st.com-lj.classic.S} on com-LiveJournal and \val{gr.st.com-orkut.classic.S} on com-Orkut, efficiencies of
\val{gr.st.com-lj.classic.E} and \val{gr.st.com-orkut.classic.E}, and on one full node the efficiencies are
\val{gr.st.com-lj.classic.Enode} and \val{gr.st.com-orkut.classic.Enode}. The one-core runs on com-Orkut were solved
once (Section~\ref{sec:timing}).

\begin{table}[H]\centering\small
%% generated by scripts/report_macros.py; do not edit
\begin{tabular}{llrr}
\toprule
dataset & CG & $S$ & $E$ \\
\midrule
\multicolumn{4}{c}{\pendingrow{vanda graph-strong: 32 of 40 runs recorded}} \\
\bottomrule
\end{tabular}

\caption{Graph strong scaling on \val{cl.gr.strong}: time per system $T$ on one core and on the largest core count,
speedup $S$ and efficiency $E$.}\label{tab:graphstrong}
\end{table}

\figslot{fig_graph_comm}{0.98\linewidth}{Time per iteration of graph CG on \val{cl.gr.strong} in local products and
vector updates, ghost exchange and reductions, classic CG.\val{spread.cap}}

At the largest core count the ghost exchange takes \val{gr.comm.com-lj.halo} of the time on com-LiveJournal and
\val{gr.comm.com-orkut.halo} on com-Orkut, and the reductions \val{gr.comm.com-lj.red} and \val{gr.comm.com-orkut.red}.
A rank there exchanges with up to \val{gr.comm.com-orkut.nbrs} others on com-Orkut and holds up to
\val{gr.comm.com-orkut.ghosts} ghosts, because a graph with hubs connects almost every rank to almost every other,
unlike the grid, whose halo is a thin surface.

On two nodes the graph runs are slower than on one. With classic CG, two nodes take
\val{pgs.com-orkut.classic.slowdown} times as long as one on com-Orkut (\val{pgs.com-lj.classic.slowdown} on
com-LiveJournal), and the ghost exchange and the reductions take \val{pgs.com-orkut.classic.2node.halo} and
\val{pgs.com-orkut.classic.2node.red} of the time. Pipelined CG is the faster method there, at
\val{pgs.com-orkut.pratio.2node} of the classic time on com-Orkut and \val{pgs.com-lj.pratio.2node} on com-LiveJournal,
because hiding a reduction behind the product is worth more than its extra vector traffic once each reduction waits for
the network. It also needs \val{pgs.com-orkut.itratio} times as many iterations as classic CG on com-Orkut
(\val{pgs.com-lj.itratio} on com-LiveJournal), a cost of its less exact recurrences.

\subsection{Partitions and exchanges}
\begin{table}[H]\centering\small
%% generated by scripts/report_macros.py; do not edit
\begin{tabular}{lrlrrrrr}
\toprule
dataset & cores & order, partition & edges cut & max ghosts & max neighbours & nnz imbalance & time/solve (s) \\
\midrule
com-LiveJournal & 288 & natural, rows & 82.1\% & 442{,}543 & 287 & 8.52 & 6.91 \\
com-LiveJournal & 288 & natural, nnz & 84.9\% & 163{,}302 & 287 & 1.00 & 5.55 \\
com-LiveJournal & 288 & bfs, nnz & 89.5\% & 183{,}242 & 254 & 1.01 & 179 \\
com-Orkut & 72 & natural, rows & 87.8\% & 2{,}002{,}066 & 71 & 3.46 & 42.7 \\
com-Orkut & 72 & natural, nnz & 89.2\% & 1{,}210{,}552 & 71 & 1.00 & 27.6 \\
com-Orkut & 72 & bfs, nnz & 91.4\% & 1{,}352{,}194 & 71 & 1.00 & 27.9 \\
\bottomrule
\end{tabular}

\caption{The three partitions at the largest core count of the method comparison on \val{cl.gr.variants}: share of
edges cut, ghosts and neighbours of the rank with the most, nonzero imbalance, and time per system with the
neighbourhood collective and classic CG.}\label{tab:graphpart}
\end{table}

\begin{table}[H]\centering\footnotesize
%% generated by scripts/report_macros.py; do not edit
\begin{tabular}{lllrrrrr}
\toprule
dataset & order, partition & layout & cut & ghosts & nbrs & imbal. & s/system \\
\midrule
com-LiveJournal & natural, rows & 72$\times$1 & 76.6\% & 974{,}437 & 71 & 5.85 & 11.0 \\
com-LiveJournal & natural, nnz & 72$\times$1 & 81.4\% & 484{,}388 & 71 & 1.00 & 5.40 \\
com-LiveJournal & bfs, nnz & 72$\times$1 & 86.9\% & 520{,}130 & 64 & 1.00 & 5.66 \\
com-LiveJournal & natural, nnz & 8$\times$9 & 63.8\% & 1{,}619{,}401 & 7 & 1.00 & 5.00 \\
com-LiveJournal & natural, nnz & 2$\times$36 & 27.2\% & 2{,}056{,}701 & 1 & 1.00 & 4.91 \\
com-Orkut & natural, rows & 72$\times$1 & 87.8\% & 2{,}002{,}066 & 71 & 3.46 & 42.7 \\
com-Orkut & natural, nnz & 72$\times$1 & 89.2\% & 1{,}210{,}552 & 71 & 1.00 & 27.6 \\
com-Orkut & bfs, nnz & 72$\times$1 & 91.4\% & 1{,}352{,}194 & 71 & 1.00 & 27.9 \\
com-Orkut & natural, nnz & 8$\times$9 & 69.4\% & 2{,}845{,}052 & 7 & 1.00 & 16.0 \\
com-Orkut & natural, nnz & 2$\times$36 & 32.7\% & 1{,}920{,}269 & 1 & 1.00 & 14.8 \\
\bottomrule
\end{tabular}

\caption{The partitions and the ranks $\times$ threads layouts on one node (\val{node.cpn} cores) of
\val{cl.gr.variants}: share of edges cut, ghosts and neighbours of the rank with the most, nonzero imbalance and time per
system, with the neighbourhood collective and classic CG. The hybrid layouts use the natural order and equal
nonzeros.}\label{tab:graphpart72}
\end{table}

\figslot{fig_graph_variants}{0.98\linewidth}{Time per system for every partition, ghost exchange and CG variant at
\val{gr.var.cores} cores on \val{cl.gr.variants}, with the ranks $\times$ threads layouts in the bottom rows. In the
row labels, rows and nnz are the partitions by equal rows and equal nonzeros, bfs is the breadth-first order, and
alltoallv, neighbor and overlap are the all-to-all exchange, the neighbourhood collective and its non-blocking form.}

On com-Orkut at \val{gr.var.com-orkut.top} cores the fastest combination is \val{gr.var.com-orkut.best} and the slowest
takes \val{gr.var.com-orkut.spread} times as long, and on com-LiveJournal the fastest is \val{gr.var.com-lj.best} and the
slowest takes \val{gr.var.com-lj.spread} times as long. The natural order cuts \val{gr.cut.com-orkut.natural.nnz} of com-Orkut's edges
and the breadth-first order \val{gr.cut.com-orkut.bfs.nnz}, and equal rows leave a nonzero imbalance of
\val{gr.imb.com-orkut.natural.rows} against \val{gr.imb.com-orkut.natural.nnz} for equal nonzeros.

\revised{graph-parallel}%
On one node the choice of partition matters most (Table~\ref{tab:graphpart72} and Figure~\ref{fig:fig_graph_variants}).
Equal rows give the busiest rank \val{gv.com-lj.node.imb.rows} times the mean number of nonzeros on com-LiveJournal
(\val{gv.com-orkut.node.imb.rows} on com-Orkut), because a contiguous range of vertices can hold several hubs, and the
solve then takes \val{gv.com-lj.node.rows} times as long as with equal nonzeros (\val{gv.com-orkut.node.rows} on
com-Orkut). The breadth-first numbering does not help these two dense graphs: it cuts \val{gv.com-lj.node.cut.bfs} of
com-LiveJournal's edges against \val{gv.com-lj.node.cut.nat} for the natural order, and the solve takes
\val{gv.com-lj.node.bfs} times as long. With \val{node.cpn} ranks almost every rank exchanges with every other, up to
\val{gv.com-orkut.node.nbrs.nat} neighbours on com-Orkut, and most edges are cut whatever the order, because these graphs
have no narrow cut along any one-dimensional order. Among the exchanges the neighbourhood collective is the fastest, the
all-to-all over every rank takes \val{gv.com-lj.node.a2a} times as long, and the overlapped form
\val{gv.com-lj.node.ovl} times, so the overlap did not hide the exchange. The likely reason is that the MPI library
does not move the messages while the rank computes. We did not test this, and a run with an asynchronous progress
thread would settle it. Pipelined CG is also slower on one node, taking \val{gv.com-lj.node.pipe} times as long on
com-LiveJournal and \val{gv.com-orkut.node.pipe} times on com-Orkut, less than on the grid (Section~\ref{sec:mpi}).

The largest gain on one node comes from fewer ranks with more threads. Two ranks of \val{node.half} threads each reduce
com-Orkut's edge cut from \val{gv.com-orkut.node.cut.nat} to \val{gv.com-orkut.node.cut.hyb} and its neighbours per rank
to \val{gv.com-orkut.node.nbrs.hyb}, and solve \val{gv.com-orkut.node.hybrid} times as fast (\val{gv.com-lj.node.hybrid}
times on com-LiveJournal). On a graph the threads of a rank share its ghosts while more ranks multiply them, whereas on
the grid a rank's halo is a thin surface whose size depends little on how many ranks share the node.

\interpslot{graph-parallel-top}{From the runs at \val{topcores} cores: whether the hybrid layout and pipelined CG together recover
scaling at \val{topcores} cores, and which method we would use there.}

\subsection{Conclusions about the data}
\revised{graph-conclusions}%
The runs say three things about these networks. First, in most of them a few labelled members identify a community:
with \val{gr.s} seeds per community, com-Amazon, com-LiveJournal, com-Youtube and com-Orkut are labelled
\val{gr.lift.range} times as accurately as the majority baseline (Table~\ref{tab:graphacc}). Second, the communities are
not equally easy to recover from the links: those of com-DBLP and email-Eu-core are the hardest, and on them more seeds
help (Section~\ref{sec:graphacc}). Third, how sparse a network is decides how robust the labelling is. The three
sparse com- graphs, com-Amazon, com-DBLP and com-Youtube (Table~\ref{tab:graphdata}), depend on few paths and lose most
of their accuracy to unreached vertices when thinned, while the densest, com-Orkut, loses much less and mostly to wrong
labels. For labelling by diffusion a missing edge therefore matters most in a network that has few, and the cost of the
method falls in the same direction, because denser graphs have larger maximum degrees, \val{gw.com-orkut.dmax} on
com-Orkut, and need more iterations, \val{gw.com-orkut.it} per system against \val{gw.com-amazon.it} on com-Amazon.

% ======================================================================================================
\section{Discussion}\label{sec:discussion}
\req{E8}
The parts of the solver are parallelised with different techniques because they have different dependences. The
stencil and the vector updates touch each point independently, so they are split by domain decomposition across ranks
and by loop sharing across threads, and they need no communication except the halo exchange. The inner products need
every point, so each rank sums its own part with its threads and the ranks then combine the partial sums in a
reduction, blocking in classic CG and non-blocking in pipelined CG. The halo exchange is point-to-point between
neighbouring ranks only, and its volume depends on the decomposition. On a GPU the same loops run as device kernels, and
the halo also travels between device and host.

Which part needs the most communication depends on the scale. By volume it is the halo exchange: at \val{topcores} ranks
a 3D block sends \val{geo.3d.cube.halo} values per exchange, while a reduction sends one or two numbers. By the number of
waits it is the reductions, two per iteration in classic CG, each involving every rank. In the measured 3D runs at
\val{meas.top} cores on \val{cl.strong} the halo exchange took \val{cm.3.halo} and the reductions \val{cm.3.red} of the
time.

\interpslot{discussion}{State the conclusions: which decomposition, halo exchange, CG variant and rank-thread layout we
would use at which scale, and what limited scaling in our runs, keeping to what the tables show.}

\emph{Limitations.} The timings use a fixed number of iterations, so they describe the cost of one iteration, and the
total cost to a tolerance, which also grows with $n$ through the iteration count, was measured only in serial. Reduction
times include waiting for the slowest rank, so they mix network latency with load imbalance. The OpenMP GPU offload
version was developed on a laptop, where the offload directives run on the CPU. It first ran on a GPU in the Vanda job
that validated it on the GPU node and then timed it, so it was not tuned for the GPU.

% ======================================================================================================
\section{Declaration of AI use}
\req{E15}
We used Claude Code (Anthropic, Claude Opus 5.5) as a coordinating session and a delegated builder session to plan the
work, write and test the C++ solver and its checks, run the experiments on the NUS HPC clusters, and draft and revise this
report, including a paragraph-by-paragraph language revision done by the AI under our direction. We directed the work and
are responsible for the content and quality of this submission.

% ======================================================================================================
\section{Submission contents and how to run}
\req{E16}\req{E17}
The source code is listed in full in Appendix~\ref{app:code}, together with the README, which explains how to build,
validate and run it. The program builds with CMake or with one \code{mpicxx} line per binary (\code{scripts/build.sh}),
\code{p3cg -{}-help} lists its options, and \code{checks/validate.py} runs the validation tests of
Section~\ref{sec:validation}. Every table and figure was generated on \val{generated_at} by scripts in the code package
from the records of \val{env.jobs} batch jobs, which ran configurations out of \val{plan.vanda.cells} planned in
advance, and these records identify each program by the start of its SHA-256 hash, \code{\val{val.sha}} for the program
validated in run \val{val.stamp} and \code{\val{val.cl.sha}} for the timed runs. The SNAP files are downloaded on the
cluster and are not part of the package, and no log files or terminal output are included.


\begin{thebibliography}{9}
\bibitem{hs} M.\,R.~Hestenes and E.~Stiefel. Methods of conjugate gradients for solving linear systems.
\emph{Journal of Research of the National Bureau of Standards} 49(6), 1952.
\bibitem{gv} P.~Ghysels and W.~Vanroose. Hiding global synchronization latency in the preconditioned conjugate
gradient algorithm. \emph{Parallel Computing} 40(7), 2014.
\bibitem{yl} J.~Yang and J.~Leskovec. Defining and evaluating network communities based on ground-truth.
\emph{IEEE International Conference on Data Mining (ICDM)}, 2012.
\bibitem{l12} A.~Chen. DSA5208 Lecture 12: The design of parallel algorithms. NUS, 2026.
\end{thebibliography}

\appendix
\section{Source code}\label{app:code}
\req{E16}
The files below are the complete solver, its build and cluster scripts, the validation programs and the README. The
scripts that set up the experiments and produce the tables and figures are part of the code package but are not listed
here.
\makeatletter
\@for\pf:=\appendixfiles\do{%
  \subsection*{\texttt{\detokenize\expandafter{\pf}}}%
  \lstinputlisting{generated/listing/\pf}}
\makeatother

\end{document}
